% ============================================================================
%  preamble/bib.tex  --  citations and reference list
% ----------------------------------------------------------------------------
%  Numeric citations. \cite{key} does TWO things:
%
%      1. prints  [1]  where you wrote it
%      2. prints the FULL reference, small, at the bottom of THAT slide
%
%  and the same work appears again, with the same number, on the references
%  slide at the end. One command, nothing to keep in sync by hand.
%
%  In your deck:
%      % ============================================================================
%  preamble/bib.tex  --  citations and reference list
% ----------------------------------------------------------------------------
%  Numeric citations. \cite{key} does TWO things:
%
%      1. prints  [1]  where you wrote it
%      2. prints the FULL reference, small, at the bottom of THAT slide
%
%  and the same work appears again, with the same number, on the references
%  slide at the end. One command, nothing to keep in sync by hand.
%
%  In your deck:
%      % ============================================================================
%  preamble/bib.tex  --  citations and reference list
% ----------------------------------------------------------------------------
%  Numeric citations. \cite{key} does TWO things:
%
%      1. prints  [1]  where you wrote it
%      2. prints the FULL reference, small, at the bottom of THAT slide
%
%  and the same work appears again, with the same number, on the references
%  slide at the end. One command, nothing to keep in sync by hand.
%
%  In your deck:
%      % ============================================================================
%  preamble/bib.tex  --  citations and reference list
% ----------------------------------------------------------------------------
%  Numeric citations. \cite{key} does TWO things:
%
%      1. prints  [1]  where you wrote it
%      2. prints the FULL reference, small, at the bottom of THAT slide
%
%  and the same work appears again, with the same number, on the references
%  slide at the end. One command, nothing to keep in sync by hand.
%
%  In your deck:
%      \input{bib}
%      \addbibresource{refs.bib}
%      ...
%      Shown by \cite{knuth1984}.   ->  "Shown by [1]."  +  [1] at slide bottom
%      \citequiet{knuth1984}        ->  [1] only, nothing at the bottom
%      \cite*{knuth1984}            ->  same as \citequiet, shorter to type
%      \bibliographyframe           ->  the references slide at the end
%
%  \cite takes biblatex's usual optional arguments:
%      \cite[p.~4]{key}          \cite[see][p.~4]{key}
%
%  ---------------------------------------------------------------------------
%  KNOBS -- put any of these in main.tex AFTER \input{bib}
%  ---------------------------------------------------------------------------
%      \deckcitefootsize{\scriptsize}  larger bottom text   (default \tiny)
%      \deckcitefootrule{0pt}          drop the hairline above the block
%      \deckcitefootmax{6}             raise the crowding warning threshold
%      \deckcitefootoff                stop adding bottom references
%      \deckcitefooton                 turn them back on
%
%  \deckcitefootoff and \deckcitefooton also work INSIDE a frame, where they
%  apply to that frame only -- useful on a related-work slide with six cites.
%
%  ---------------------------------------------------------------------------
%  WHERE NOT TO PUT \cite
%  ---------------------------------------------------------------------------
%  Not in \frametitle, \section or \caption. Those are moving arguments: the
%  text is written out to the .toc/.aux and typeset again elsewhere, and a
%  footnote cannot survive that. Use \citequiet there.
% ============================================================================

\usepackage[
  backend     = biber,
  style       = numeric,   % [1] in the text, [1] in the reference list
  sorting     = none,      % list ordered by first citation, so [1] really is
                           % the first work you cite -- numbers read in order
  maxbibnames = 6,
  doi         = false,     % slides rarely need these; they eat space
  isbn        = false,
  url         = false,
  eprint      = false,
]{biblatex}

\usepackage{etoolbox}      % \forcsvlist -- biblatex loads it too, this line
                           % is here so the file does not rely on that

% ---------------------------------------------------------------------------
%  THE REFERENCES SLIDE
% ---------------------------------------------------------------------------
%  beamer replaces biblatex's entry label with its own "bibliography item"
%  template. Left empty it drops beamer's generic article icon -- but it also
%  silently deletes the [1] that numeric style depends on. \insertbiblabel
%  puts the number back.
\setbeamertemplate{bibliography item}{\insertbiblabel}
\renewcommand*{\bibfont}{\scriptsize}

% Keep entries from being hyphenated into a mess in a narrow column
\appto{\bibsetup}{\raggedright}

%  A frame that breaks across pages is titled "References I", "References II"
%  and so on. Numbering a list that turned out to fit on one slide just looks
%  like a mistake, so the numeral starts from the second page.
\setbeamertemplate{frametitle continuation}[from second]

% ---------------------------------------------------------------------------
%  A references slide that splits across pages automatically if it is long.
%  Use \bibliographyframe in your deck instead of writing this out by hand.
% ---------------------------------------------------------------------------
\newcommand{\bibliographyframe}[1][References]{%
  \begin{frame}[allowframebreaks]{#1}
    \printbibliography[heading=none]
  \end{frame}%
}

% ===========================================================================
%  THE PER-SLIDE REFERENCE BLOCK
% ===========================================================================
\makeatletter

% --- Knobs -----------------------------------------------------------------
\newcommand*{\deck@citefootsize}{\tiny}
\newcommand*{\deckcitefootsize}[1]{\renewcommand*{\deck@citefootsize}{#1}}

\newlength{\deck@citefootrule}
\setlength{\deck@citefootrule}{0.4pt}
\newcommand*{\deckcitefootrule}[1]{\setlength{\deck@citefootrule}{#1}}

\newcount\deck@citefootmax  \deck@citefootmax=4
\newcommand*{\deckcitefootmax}[1]{\deck@citefootmax=#1\relax}

\newif\ifdeck@citefoot  \deck@citefoottrue
\newcommand*{\deckcitefooton}{\deck@citefoottrue}
\newcommand*{\deckcitefootoff}{\deck@citefootfalse}

% --- The number, on its own ------------------------------------------------
%  Prints just [1] for a key, so the bottom line can be labelled with the
%  same number the reader just saw in the body text.
\DeclareCiteCommand{\deckcitenum}[\mkbibbrackets]
  {}
  {\printfield{labelnumber}}
  {\multicitedelim}
  {}

% --- The hairline above the block ------------------------------------------
%  This is LaTeX's footnote rule, so an ordinary \footnote gets the same line.
%  \deckcitefootrule{0pt} makes it disappear without disturbing the spacing.
\renewcommand{\footnoterule}{%
  \kern-3pt%
  \hbox{\color{BrandRule}%
        \vrule width 0.35\textwidth height \deck@citefootrule depth 0pt}%
  \kern2.6pt%
}

% --- One bottom line -------------------------------------------------------
%  An unmarked footnote: beamer already places footnotes at the bottom of the
%  frame, above the footer, and keeps them working under [allowframebreaks].
%  The mark is emptied and the footnote counter rolled back, because the line
%  carries its own [n] instead.
%
%  The footnote template is overridden LOCALLY. beamer's default reserves
%  1.8em for the mark; with no mark that is 1.8em of dead indent on every
%  line. Ordinary footnotes elsewhere in the deck are untouched.
%
%  The four \setbeamerfont lines are not optional. \fullcite goes through
%  beamer's bibliography drivers, which apply "bibliography entry ..." fonts --
%  absolute sizes that would override \deckcitefootsize and leave the [n]
%  label tiny next to a much larger reference. Unstarred \setbeamerfont
%  changes only the size and leaves series and shape alone, and all four are
%  local to the group, so the references slide at the end is unaffected.
\newcommand{\deck@citefootline}[1]{%
  \begingroup
    \setbeamerfont{bibliography entry author}  {size=\deck@citefootsize}%
    \setbeamerfont{bibliography entry title}   {size=\deck@citefootsize}%
    \setbeamerfont{bibliography entry location}{size=\deck@citefootsize}%
    \setbeamerfont{bibliography entry note}    {size=\deck@citefootsize}%
    \setbeamertemplate{footnote}{%
      \parindent0pt\noindent\raggedright
      \insertfootnotetext\par
    }%
    \renewcommand{\thefootnote}{}%
    \footnote{{\deck@citefootsize\deckcitenum{#1}\ \fullcite{#1}}}%
    \addtocounter{footnote}{-1}%
  \endgroup
}

% --- Crowding warning ------------------------------------------------------
%  Counts bottom references on the current slide and says so in the log once
%  the slide carries more than \deckcitefootmax of them. A crushed slide
%  should never be a surprise you only discover in the PDF.
\newcount\deck@citefootn
\def\deck@citefootwhere{}
\newcommand{\deck@citefootcount}{%
  \edef\deck@citefoothere{\the\c@page.\the\beamer@slideinframe}%
  \ifx\deck@citefoothere\deck@citefootwhere
    \global\advance\deck@citefootn\@ne
  \else
    \global\deck@citefootn\@ne
    \global\let\deck@citefootwhere\deck@citefoothere
  \fi
  \ifnum\deck@citefootn>\deck@citefootmax\relax
    \PackageWarning{deck-bib}{%
      \the\deck@citefootn\space bottom references on slide
      \the\c@page\space -- more than \the\deck@citefootmax.\MessageBreak
      They will crowd the slide. Use \string\citequiet\space for the
      extras,\MessageBreak
      or \string\deckcitefootoff\space at the top of that frame%
    }%
  \fi
}

% --- Dedup, per slide ------------------------------------------------------
%  Cite the same work twice on one slide and it should appear once at the
%  bottom; cite it again on a later slide and it should appear there too.
%
%  The marker is keyed on the physical page AND on which overlay of the frame
%  is being built. beamer typesets a frame body once per overlay, so a flag
%  reset by a begin-frame hook would fire once and then suppress the block on
%  overlays 2..n -- the reference would vanish the moment you pressed the
%  clicker. Keying on the slide being built needs no hook at all.
\newcommand{\deck@citefootone}[1]{%
  \ifcsname deck@seen@\the\c@page @\the\beamer@slideinframe @#1\endcsname
  \else
    \expandafter\gdef
      \csname deck@seen@\the\c@page @\the\beamer@slideinframe @#1\endcsname{}%
    \deck@citefootcount
    \deck@citefootline{#1}%
  \fi
}

%  \cite{a,b} is legal, so split the list. \forcsvlist trims stray spaces.
\newcommand{\deck@citefoot}[1]{%
  \ifdeck@citefoot
    \forcsvlist{\deck@citefootone}{#1}%
  \fi
}

% --- \cite does both -------------------------------------------------------
%  Wrapped, not reimplemented, so biblatex keeps full control of what [1]
%  actually looks like.
%
%  Deferred to \AtBeginDocument because hyperref -- which beamer loads for us
%  -- patches \cite at the end of the preamble. Grabbing it any earlier would
%  wrap a version that is about to be replaced.
%  biblatex reads ONE optional argument as the postnote and TWO as prenote +
%  postnote, so the wrapper has to hand back exactly as many as it was given.
%  Empty defaults are used rather than -NoValue- markers: an empty argument is
%  something \detokenize can test for certainty, and there is nothing to pass
%  through a second command that might mistake the marker for real text.
\newcommand{\deck@ifempty}[1]{%
  \if\relax\detokenize{#1}\relax
    \expandafter\@firstoftwo
  \else
    \expandafter\@secondoftwo
  \fi
}
\newcommand{\deck@passcite}[3]{%
  \deck@ifempty{#1}%
    {\deck@ifempty{#2}{\deck@origcite{#3}}{\deck@origcite[][#2]{#3}}}%
    {\deck@ifempty{#2}{\deck@origcite[#1]{#3}}{\deck@origcite[#1][#2]{#3}}}%
}

\NewDocumentCommand{\citequiet}{O{} O{} m}{\deck@passcite{#1}{#2}{#3}}

\AtBeginDocument{%
  \let\deck@origcite\cite
  \RenewDocumentCommand{\cite}{s O{} O{} m}{%
    \deck@passcite{#2}{#3}{#4}%
    \IfBooleanF{#1}{\deck@citefoot{#4}}%
  }%
}

\makeatother

%      \addbibresource{refs.bib}
%      ...
%      Shown by \cite{knuth1984}.   ->  "Shown by [1]."  +  [1] at slide bottom
%      \citequiet{knuth1984}        ->  [1] only, nothing at the bottom
%      \cite*{knuth1984}            ->  same as \citequiet, shorter to type
%      \bibliographyframe           ->  the references slide at the end
%
%  \cite takes biblatex's usual optional arguments:
%      \cite[p.~4]{key}          \cite[see][p.~4]{key}
%
%  ---------------------------------------------------------------------------
%  KNOBS -- put any of these in main.tex AFTER % ============================================================================
%  preamble/bib.tex  --  citations and reference list
% ----------------------------------------------------------------------------
%  Numeric citations. \cite{key} does TWO things:
%
%      1. prints  [1]  where you wrote it
%      2. prints the FULL reference, small, at the bottom of THAT slide
%
%  and the same work appears again, with the same number, on the references
%  slide at the end. One command, nothing to keep in sync by hand.
%
%  In your deck:
%      \input{bib}
%      \addbibresource{refs.bib}
%      ...
%      Shown by \cite{knuth1984}.   ->  "Shown by [1]."  +  [1] at slide bottom
%      \citequiet{knuth1984}        ->  [1] only, nothing at the bottom
%      \cite*{knuth1984}            ->  same as \citequiet, shorter to type
%      \bibliographyframe           ->  the references slide at the end
%
%  \cite takes biblatex's usual optional arguments:
%      \cite[p.~4]{key}          \cite[see][p.~4]{key}
%
%  ---------------------------------------------------------------------------
%  KNOBS -- put any of these in main.tex AFTER \input{bib}
%  ---------------------------------------------------------------------------
%      \deckcitefootsize{\scriptsize}  larger bottom text   (default \tiny)
%      \deckcitefootrule{0pt}          drop the hairline above the block
%      \deckcitefootmax{6}             raise the crowding warning threshold
%      \deckcitefootoff                stop adding bottom references
%      \deckcitefooton                 turn them back on
%
%  \deckcitefootoff and \deckcitefooton also work INSIDE a frame, where they
%  apply to that frame only -- useful on a related-work slide with six cites.
%
%  ---------------------------------------------------------------------------
%  WHERE NOT TO PUT \cite
%  ---------------------------------------------------------------------------
%  Not in \frametitle, \section or \caption. Those are moving arguments: the
%  text is written out to the .toc/.aux and typeset again elsewhere, and a
%  footnote cannot survive that. Use \citequiet there.
% ============================================================================

\usepackage[
  backend     = biber,
  style       = numeric,   % [1] in the text, [1] in the reference list
  sorting     = none,      % list ordered by first citation, so [1] really is
                           % the first work you cite -- numbers read in order
  maxbibnames = 6,
  doi         = false,     % slides rarely need these; they eat space
  isbn        = false,
  url         = false,
  eprint      = false,
]{biblatex}

\usepackage{etoolbox}      % \forcsvlist -- biblatex loads it too, this line
                           % is here so the file does not rely on that

% ---------------------------------------------------------------------------
%  THE REFERENCES SLIDE
% ---------------------------------------------------------------------------
%  beamer replaces biblatex's entry label with its own "bibliography item"
%  template. Left empty it drops beamer's generic article icon -- but it also
%  silently deletes the [1] that numeric style depends on. \insertbiblabel
%  puts the number back.
\setbeamertemplate{bibliography item}{\insertbiblabel}
\renewcommand*{\bibfont}{\scriptsize}

% Keep entries from being hyphenated into a mess in a narrow column
\appto{\bibsetup}{\raggedright}

%  A frame that breaks across pages is titled "References I", "References II"
%  and so on. Numbering a list that turned out to fit on one slide just looks
%  like a mistake, so the numeral starts from the second page.
\setbeamertemplate{frametitle continuation}[from second]

% ---------------------------------------------------------------------------
%  A references slide that splits across pages automatically if it is long.
%  Use \bibliographyframe in your deck instead of writing this out by hand.
% ---------------------------------------------------------------------------
\newcommand{\bibliographyframe}[1][References]{%
  \begin{frame}[allowframebreaks]{#1}
    \printbibliography[heading=none]
  \end{frame}%
}

% ===========================================================================
%  THE PER-SLIDE REFERENCE BLOCK
% ===========================================================================
\makeatletter

% --- Knobs -----------------------------------------------------------------
\newcommand*{\deck@citefootsize}{\tiny}
\newcommand*{\deckcitefootsize}[1]{\renewcommand*{\deck@citefootsize}{#1}}

\newlength{\deck@citefootrule}
\setlength{\deck@citefootrule}{0.4pt}
\newcommand*{\deckcitefootrule}[1]{\setlength{\deck@citefootrule}{#1}}

\newcount\deck@citefootmax  \deck@citefootmax=4
\newcommand*{\deckcitefootmax}[1]{\deck@citefootmax=#1\relax}

\newif\ifdeck@citefoot  \deck@citefoottrue
\newcommand*{\deckcitefooton}{\deck@citefoottrue}
\newcommand*{\deckcitefootoff}{\deck@citefootfalse}

% --- The number, on its own ------------------------------------------------
%  Prints just [1] for a key, so the bottom line can be labelled with the
%  same number the reader just saw in the body text.
\DeclareCiteCommand{\deckcitenum}[\mkbibbrackets]
  {}
  {\printfield{labelnumber}}
  {\multicitedelim}
  {}

% --- The hairline above the block ------------------------------------------
%  This is LaTeX's footnote rule, so an ordinary \footnote gets the same line.
%  \deckcitefootrule{0pt} makes it disappear without disturbing the spacing.
\renewcommand{\footnoterule}{%
  \kern-3pt%
  \hbox{\color{BrandRule}%
        \vrule width 0.35\textwidth height \deck@citefootrule depth 0pt}%
  \kern2.6pt%
}

% --- One bottom line -------------------------------------------------------
%  An unmarked footnote: beamer already places footnotes at the bottom of the
%  frame, above the footer, and keeps them working under [allowframebreaks].
%  The mark is emptied and the footnote counter rolled back, because the line
%  carries its own [n] instead.
%
%  The footnote template is overridden LOCALLY. beamer's default reserves
%  1.8em for the mark; with no mark that is 1.8em of dead indent on every
%  line. Ordinary footnotes elsewhere in the deck are untouched.
%
%  The four \setbeamerfont lines are not optional. \fullcite goes through
%  beamer's bibliography drivers, which apply "bibliography entry ..." fonts --
%  absolute sizes that would override \deckcitefootsize and leave the [n]
%  label tiny next to a much larger reference. Unstarred \setbeamerfont
%  changes only the size and leaves series and shape alone, and all four are
%  local to the group, so the references slide at the end is unaffected.
\newcommand{\deck@citefootline}[1]{%
  \begingroup
    \setbeamerfont{bibliography entry author}  {size=\deck@citefootsize}%
    \setbeamerfont{bibliography entry title}   {size=\deck@citefootsize}%
    \setbeamerfont{bibliography entry location}{size=\deck@citefootsize}%
    \setbeamerfont{bibliography entry note}    {size=\deck@citefootsize}%
    \setbeamertemplate{footnote}{%
      \parindent0pt\noindent\raggedright
      \insertfootnotetext\par
    }%
    \renewcommand{\thefootnote}{}%
    \footnote{{\deck@citefootsize\deckcitenum{#1}\ \fullcite{#1}}}%
    \addtocounter{footnote}{-1}%
  \endgroup
}

% --- Crowding warning ------------------------------------------------------
%  Counts bottom references on the current slide and says so in the log once
%  the slide carries more than \deckcitefootmax of them. A crushed slide
%  should never be a surprise you only discover in the PDF.
\newcount\deck@citefootn
\def\deck@citefootwhere{}
\newcommand{\deck@citefootcount}{%
  \edef\deck@citefoothere{\the\c@page.\the\beamer@slideinframe}%
  \ifx\deck@citefoothere\deck@citefootwhere
    \global\advance\deck@citefootn\@ne
  \else
    \global\deck@citefootn\@ne
    \global\let\deck@citefootwhere\deck@citefoothere
  \fi
  \ifnum\deck@citefootn>\deck@citefootmax\relax
    \PackageWarning{deck-bib}{%
      \the\deck@citefootn\space bottom references on slide
      \the\c@page\space -- more than \the\deck@citefootmax.\MessageBreak
      They will crowd the slide. Use \string\citequiet\space for the
      extras,\MessageBreak
      or \string\deckcitefootoff\space at the top of that frame%
    }%
  \fi
}

% --- Dedup, per slide ------------------------------------------------------
%  Cite the same work twice on one slide and it should appear once at the
%  bottom; cite it again on a later slide and it should appear there too.
%
%  The marker is keyed on the physical page AND on which overlay of the frame
%  is being built. beamer typesets a frame body once per overlay, so a flag
%  reset by a begin-frame hook would fire once and then suppress the block on
%  overlays 2..n -- the reference would vanish the moment you pressed the
%  clicker. Keying on the slide being built needs no hook at all.
\newcommand{\deck@citefootone}[1]{%
  \ifcsname deck@seen@\the\c@page @\the\beamer@slideinframe @#1\endcsname
  \else
    \expandafter\gdef
      \csname deck@seen@\the\c@page @\the\beamer@slideinframe @#1\endcsname{}%
    \deck@citefootcount
    \deck@citefootline{#1}%
  \fi
}

%  \cite{a,b} is legal, so split the list. \forcsvlist trims stray spaces.
\newcommand{\deck@citefoot}[1]{%
  \ifdeck@citefoot
    \forcsvlist{\deck@citefootone}{#1}%
  \fi
}

% --- \cite does both -------------------------------------------------------
%  Wrapped, not reimplemented, so biblatex keeps full control of what [1]
%  actually looks like.
%
%  Deferred to \AtBeginDocument because hyperref -- which beamer loads for us
%  -- patches \cite at the end of the preamble. Grabbing it any earlier would
%  wrap a version that is about to be replaced.
%  biblatex reads ONE optional argument as the postnote and TWO as prenote +
%  postnote, so the wrapper has to hand back exactly as many as it was given.
%  Empty defaults are used rather than -NoValue- markers: an empty argument is
%  something \detokenize can test for certainty, and there is nothing to pass
%  through a second command that might mistake the marker for real text.
\newcommand{\deck@ifempty}[1]{%
  \if\relax\detokenize{#1}\relax
    \expandafter\@firstoftwo
  \else
    \expandafter\@secondoftwo
  \fi
}
\newcommand{\deck@passcite}[3]{%
  \deck@ifempty{#1}%
    {\deck@ifempty{#2}{\deck@origcite{#3}}{\deck@origcite[][#2]{#3}}}%
    {\deck@ifempty{#2}{\deck@origcite[#1]{#3}}{\deck@origcite[#1][#2]{#3}}}%
}

\NewDocumentCommand{\citequiet}{O{} O{} m}{\deck@passcite{#1}{#2}{#3}}

\AtBeginDocument{%
  \let\deck@origcite\cite
  \RenewDocumentCommand{\cite}{s O{} O{} m}{%
    \deck@passcite{#2}{#3}{#4}%
    \IfBooleanF{#1}{\deck@citefoot{#4}}%
  }%
}

\makeatother

%  ---------------------------------------------------------------------------
%      \deckcitefootsize{\scriptsize}  larger bottom text   (default \tiny)
%      \deckcitefootrule{0pt}          drop the hairline above the block
%      \deckcitefootmax{6}             raise the crowding warning threshold
%      \deckcitefootoff                stop adding bottom references
%      \deckcitefooton                 turn them back on
%
%  \deckcitefootoff and \deckcitefooton also work INSIDE a frame, where they
%  apply to that frame only -- useful on a related-work slide with six cites.
%
%  ---------------------------------------------------------------------------
%  WHERE NOT TO PUT \cite
%  ---------------------------------------------------------------------------
%  Not in \frametitle, \section or \caption. Those are moving arguments: the
%  text is written out to the .toc/.aux and typeset again elsewhere, and a
%  footnote cannot survive that. Use \citequiet there.
% ============================================================================

\usepackage[
  backend     = biber,
  style       = numeric,   % [1] in the text, [1] in the reference list
  sorting     = none,      % list ordered by first citation, so [1] really is
                           % the first work you cite -- numbers read in order
  maxbibnames = 6,
  doi         = false,     % slides rarely need these; they eat space
  isbn        = false,
  url         = false,
  eprint      = false,
]{biblatex}

\usepackage{etoolbox}      % \forcsvlist -- biblatex loads it too, this line
                           % is here so the file does not rely on that

% ---------------------------------------------------------------------------
%  THE REFERENCES SLIDE
% ---------------------------------------------------------------------------
%  beamer replaces biblatex's entry label with its own "bibliography item"
%  template. Left empty it drops beamer's generic article icon -- but it also
%  silently deletes the [1] that numeric style depends on. \insertbiblabel
%  puts the number back.
\setbeamertemplate{bibliography item}{\insertbiblabel}
\renewcommand*{\bibfont}{\scriptsize}

% Keep entries from being hyphenated into a mess in a narrow column
\appto{\bibsetup}{\raggedright}

%  A frame that breaks across pages is titled "References I", "References II"
%  and so on. Numbering a list that turned out to fit on one slide just looks
%  like a mistake, so the numeral starts from the second page.
\setbeamertemplate{frametitle continuation}[from second]

% ---------------------------------------------------------------------------
%  A references slide that splits across pages automatically if it is long.
%  Use \bibliographyframe in your deck instead of writing this out by hand.
% ---------------------------------------------------------------------------
\newcommand{\bibliographyframe}[1][References]{%
  \begin{frame}[allowframebreaks]{#1}
    \printbibliography[heading=none]
  \end{frame}%
}

% ===========================================================================
%  THE PER-SLIDE REFERENCE BLOCK
% ===========================================================================
\makeatletter

% --- Knobs -----------------------------------------------------------------
\newcommand*{\deck@citefootsize}{\tiny}
\newcommand*{\deckcitefootsize}[1]{\renewcommand*{\deck@citefootsize}{#1}}

\newlength{\deck@citefootrule}
\setlength{\deck@citefootrule}{0.4pt}
\newcommand*{\deckcitefootrule}[1]{\setlength{\deck@citefootrule}{#1}}

\newcount\deck@citefootmax  \deck@citefootmax=4
\newcommand*{\deckcitefootmax}[1]{\deck@citefootmax=#1\relax}

\newif\ifdeck@citefoot  \deck@citefoottrue
\newcommand*{\deckcitefooton}{\deck@citefoottrue}
\newcommand*{\deckcitefootoff}{\deck@citefootfalse}

% --- The number, on its own ------------------------------------------------
%  Prints just [1] for a key, so the bottom line can be labelled with the
%  same number the reader just saw in the body text.
\DeclareCiteCommand{\deckcitenum}[\mkbibbrackets]
  {}
  {\printfield{labelnumber}}
  {\multicitedelim}
  {}

% --- The hairline above the block ------------------------------------------
%  This is LaTeX's footnote rule, so an ordinary \footnote gets the same line.
%  \deckcitefootrule{0pt} makes it disappear without disturbing the spacing.
\renewcommand{\footnoterule}{%
  \kern-3pt%
  \hbox{\color{BrandRule}%
        \vrule width 0.35\textwidth height \deck@citefootrule depth 0pt}%
  \kern2.6pt%
}

% --- One bottom line -------------------------------------------------------
%  An unmarked footnote: beamer already places footnotes at the bottom of the
%  frame, above the footer, and keeps them working under [allowframebreaks].
%  The mark is emptied and the footnote counter rolled back, because the line
%  carries its own [n] instead.
%
%  The footnote template is overridden LOCALLY. beamer's default reserves
%  1.8em for the mark; with no mark that is 1.8em of dead indent on every
%  line. Ordinary footnotes elsewhere in the deck are untouched.
%
%  The four \setbeamerfont lines are not optional. \fullcite goes through
%  beamer's bibliography drivers, which apply "bibliography entry ..." fonts --
%  absolute sizes that would override \deckcitefootsize and leave the [n]
%  label tiny next to a much larger reference. Unstarred \setbeamerfont
%  changes only the size and leaves series and shape alone, and all four are
%  local to the group, so the references slide at the end is unaffected.
\newcommand{\deck@citefootline}[1]{%
  \begingroup
    \setbeamerfont{bibliography entry author}  {size=\deck@citefootsize}%
    \setbeamerfont{bibliography entry title}   {size=\deck@citefootsize}%
    \setbeamerfont{bibliography entry location}{size=\deck@citefootsize}%
    \setbeamerfont{bibliography entry note}    {size=\deck@citefootsize}%
    \setbeamertemplate{footnote}{%
      \parindent0pt\noindent\raggedright
      \insertfootnotetext\par
    }%
    \renewcommand{\thefootnote}{}%
    \footnote{{\deck@citefootsize\deckcitenum{#1}\ \fullcite{#1}}}%
    \addtocounter{footnote}{-1}%
  \endgroup
}

% --- Crowding warning ------------------------------------------------------
%  Counts bottom references on the current slide and says so in the log once
%  the slide carries more than \deckcitefootmax of them. A crushed slide
%  should never be a surprise you only discover in the PDF.
\newcount\deck@citefootn
\def\deck@citefootwhere{}
\newcommand{\deck@citefootcount}{%
  \edef\deck@citefoothere{\the\c@page.\the\beamer@slideinframe}%
  \ifx\deck@citefoothere\deck@citefootwhere
    \global\advance\deck@citefootn\@ne
  \else
    \global\deck@citefootn\@ne
    \global\let\deck@citefootwhere\deck@citefoothere
  \fi
  \ifnum\deck@citefootn>\deck@citefootmax\relax
    \PackageWarning{deck-bib}{%
      \the\deck@citefootn\space bottom references on slide
      \the\c@page\space -- more than \the\deck@citefootmax.\MessageBreak
      They will crowd the slide. Use \string\citequiet\space for the
      extras,\MessageBreak
      or \string\deckcitefootoff\space at the top of that frame%
    }%
  \fi
}

% --- Dedup, per slide ------------------------------------------------------
%  Cite the same work twice on one slide and it should appear once at the
%  bottom; cite it again on a later slide and it should appear there too.
%
%  The marker is keyed on the physical page AND on which overlay of the frame
%  is being built. beamer typesets a frame body once per overlay, so a flag
%  reset by a begin-frame hook would fire once and then suppress the block on
%  overlays 2..n -- the reference would vanish the moment you pressed the
%  clicker. Keying on the slide being built needs no hook at all.
\newcommand{\deck@citefootone}[1]{%
  \ifcsname deck@seen@\the\c@page @\the\beamer@slideinframe @#1\endcsname
  \else
    \expandafter\gdef
      \csname deck@seen@\the\c@page @\the\beamer@slideinframe @#1\endcsname{}%
    \deck@citefootcount
    \deck@citefootline{#1}%
  \fi
}

%  \cite{a,b} is legal, so split the list. \forcsvlist trims stray spaces.
\newcommand{\deck@citefoot}[1]{%
  \ifdeck@citefoot
    \forcsvlist{\deck@citefootone}{#1}%
  \fi
}

% --- \cite does both -------------------------------------------------------
%  Wrapped, not reimplemented, so biblatex keeps full control of what [1]
%  actually looks like.
%
%  Deferred to \AtBeginDocument because hyperref -- which beamer loads for us
%  -- patches \cite at the end of the preamble. Grabbing it any earlier would
%  wrap a version that is about to be replaced.
%  biblatex reads ONE optional argument as the postnote and TWO as prenote +
%  postnote, so the wrapper has to hand back exactly as many as it was given.
%  Empty defaults are used rather than -NoValue- markers: an empty argument is
%  something \detokenize can test for certainty, and there is nothing to pass
%  through a second command that might mistake the marker for real text.
\newcommand{\deck@ifempty}[1]{%
  \if\relax\detokenize{#1}\relax
    \expandafter\@firstoftwo
  \else
    \expandafter\@secondoftwo
  \fi
}
\newcommand{\deck@passcite}[3]{%
  \deck@ifempty{#1}%
    {\deck@ifempty{#2}{\deck@origcite{#3}}{\deck@origcite[][#2]{#3}}}%
    {\deck@ifempty{#2}{\deck@origcite[#1]{#3}}{\deck@origcite[#1][#2]{#3}}}%
}

\NewDocumentCommand{\citequiet}{O{} O{} m}{\deck@passcite{#1}{#2}{#3}}

\AtBeginDocument{%
  \let\deck@origcite\cite
  \RenewDocumentCommand{\cite}{s O{} O{} m}{%
    \deck@passcite{#2}{#3}{#4}%
    \IfBooleanF{#1}{\deck@citefoot{#4}}%
  }%
}

\makeatother

%      \addbibresource{refs.bib}
%      ...
%      Shown by \cite{knuth1984}.   ->  "Shown by [1]."  +  [1] at slide bottom
%      \citequiet{knuth1984}        ->  [1] only, nothing at the bottom
%      \cite*{knuth1984}            ->  same as \citequiet, shorter to type
%      \bibliographyframe           ->  the references slide at the end
%
%  \cite takes biblatex's usual optional arguments:
%      \cite[p.~4]{key}          \cite[see][p.~4]{key}
%
%  ---------------------------------------------------------------------------
%  KNOBS -- put any of these in main.tex AFTER % ============================================================================
%  preamble/bib.tex  --  citations and reference list
% ----------------------------------------------------------------------------
%  Numeric citations. \cite{key} does TWO things:
%
%      1. prints  [1]  where you wrote it
%      2. prints the FULL reference, small, at the bottom of THAT slide
%
%  and the same work appears again, with the same number, on the references
%  slide at the end. One command, nothing to keep in sync by hand.
%
%  In your deck:
%      % ============================================================================
%  preamble/bib.tex  --  citations and reference list
% ----------------------------------------------------------------------------
%  Numeric citations. \cite{key} does TWO things:
%
%      1. prints  [1]  where you wrote it
%      2. prints the FULL reference, small, at the bottom of THAT slide
%
%  and the same work appears again, with the same number, on the references
%  slide at the end. One command, nothing to keep in sync by hand.
%
%  In your deck:
%      \input{bib}
%      \addbibresource{refs.bib}
%      ...
%      Shown by \cite{knuth1984}.   ->  "Shown by [1]."  +  [1] at slide bottom
%      \citequiet{knuth1984}        ->  [1] only, nothing at the bottom
%      \cite*{knuth1984}            ->  same as \citequiet, shorter to type
%      \bibliographyframe           ->  the references slide at the end
%
%  \cite takes biblatex's usual optional arguments:
%      \cite[p.~4]{key}          \cite[see][p.~4]{key}
%
%  ---------------------------------------------------------------------------
%  KNOBS -- put any of these in main.tex AFTER \input{bib}
%  ---------------------------------------------------------------------------
%      \deckcitefootsize{\scriptsize}  larger bottom text   (default \tiny)
%      \deckcitefootrule{0pt}          drop the hairline above the block
%      \deckcitefootmax{6}             raise the crowding warning threshold
%      \deckcitefootoff                stop adding bottom references
%      \deckcitefooton                 turn them back on
%
%  \deckcitefootoff and \deckcitefooton also work INSIDE a frame, where they
%  apply to that frame only -- useful on a related-work slide with six cites.
%
%  ---------------------------------------------------------------------------
%  WHERE NOT TO PUT \cite
%  ---------------------------------------------------------------------------
%  Not in \frametitle, \section or \caption. Those are moving arguments: the
%  text is written out to the .toc/.aux and typeset again elsewhere, and a
%  footnote cannot survive that. Use \citequiet there.
% ============================================================================

\usepackage[
  backend     = biber,
  style       = numeric,   % [1] in the text, [1] in the reference list
  sorting     = none,      % list ordered by first citation, so [1] really is
                           % the first work you cite -- numbers read in order
  maxbibnames = 6,
  doi         = false,     % slides rarely need these; they eat space
  isbn        = false,
  url         = false,
  eprint      = false,
]{biblatex}

\usepackage{etoolbox}      % \forcsvlist -- biblatex loads it too, this line
                           % is here so the file does not rely on that

% ---------------------------------------------------------------------------
%  THE REFERENCES SLIDE
% ---------------------------------------------------------------------------
%  beamer replaces biblatex's entry label with its own "bibliography item"
%  template. Left empty it drops beamer's generic article icon -- but it also
%  silently deletes the [1] that numeric style depends on. \insertbiblabel
%  puts the number back.
\setbeamertemplate{bibliography item}{\insertbiblabel}
\renewcommand*{\bibfont}{\scriptsize}

% Keep entries from being hyphenated into a mess in a narrow column
\appto{\bibsetup}{\raggedright}

%  A frame that breaks across pages is titled "References I", "References II"
%  and so on. Numbering a list that turned out to fit on one slide just looks
%  like a mistake, so the numeral starts from the second page.
\setbeamertemplate{frametitle continuation}[from second]

% ---------------------------------------------------------------------------
%  A references slide that splits across pages automatically if it is long.
%  Use \bibliographyframe in your deck instead of writing this out by hand.
% ---------------------------------------------------------------------------
\newcommand{\bibliographyframe}[1][References]{%
  \begin{frame}[allowframebreaks]{#1}
    \printbibliography[heading=none]
  \end{frame}%
}

% ===========================================================================
%  THE PER-SLIDE REFERENCE BLOCK
% ===========================================================================
\makeatletter

% --- Knobs -----------------------------------------------------------------
\newcommand*{\deck@citefootsize}{\tiny}
\newcommand*{\deckcitefootsize}[1]{\renewcommand*{\deck@citefootsize}{#1}}

\newlength{\deck@citefootrule}
\setlength{\deck@citefootrule}{0.4pt}
\newcommand*{\deckcitefootrule}[1]{\setlength{\deck@citefootrule}{#1}}

\newcount\deck@citefootmax  \deck@citefootmax=4
\newcommand*{\deckcitefootmax}[1]{\deck@citefootmax=#1\relax}

\newif\ifdeck@citefoot  \deck@citefoottrue
\newcommand*{\deckcitefooton}{\deck@citefoottrue}
\newcommand*{\deckcitefootoff}{\deck@citefootfalse}

% --- The number, on its own ------------------------------------------------
%  Prints just [1] for a key, so the bottom line can be labelled with the
%  same number the reader just saw in the body text.
\DeclareCiteCommand{\deckcitenum}[\mkbibbrackets]
  {}
  {\printfield{labelnumber}}
  {\multicitedelim}
  {}

% --- The hairline above the block ------------------------------------------
%  This is LaTeX's footnote rule, so an ordinary \footnote gets the same line.
%  \deckcitefootrule{0pt} makes it disappear without disturbing the spacing.
\renewcommand{\footnoterule}{%
  \kern-3pt%
  \hbox{\color{BrandRule}%
        \vrule width 0.35\textwidth height \deck@citefootrule depth 0pt}%
  \kern2.6pt%
}

% --- One bottom line -------------------------------------------------------
%  An unmarked footnote: beamer already places footnotes at the bottom of the
%  frame, above the footer, and keeps them working under [allowframebreaks].
%  The mark is emptied and the footnote counter rolled back, because the line
%  carries its own [n] instead.
%
%  The footnote template is overridden LOCALLY. beamer's default reserves
%  1.8em for the mark; with no mark that is 1.8em of dead indent on every
%  line. Ordinary footnotes elsewhere in the deck are untouched.
%
%  The four \setbeamerfont lines are not optional. \fullcite goes through
%  beamer's bibliography drivers, which apply "bibliography entry ..." fonts --
%  absolute sizes that would override \deckcitefootsize and leave the [n]
%  label tiny next to a much larger reference. Unstarred \setbeamerfont
%  changes only the size and leaves series and shape alone, and all four are
%  local to the group, so the references slide at the end is unaffected.
\newcommand{\deck@citefootline}[1]{%
  \begingroup
    \setbeamerfont{bibliography entry author}  {size=\deck@citefootsize}%
    \setbeamerfont{bibliography entry title}   {size=\deck@citefootsize}%
    \setbeamerfont{bibliography entry location}{size=\deck@citefootsize}%
    \setbeamerfont{bibliography entry note}    {size=\deck@citefootsize}%
    \setbeamertemplate{footnote}{%
      \parindent0pt\noindent\raggedright
      \insertfootnotetext\par
    }%
    \renewcommand{\thefootnote}{}%
    \footnote{{\deck@citefootsize\deckcitenum{#1}\ \fullcite{#1}}}%
    \addtocounter{footnote}{-1}%
  \endgroup
}

% --- Crowding warning ------------------------------------------------------
%  Counts bottom references on the current slide and says so in the log once
%  the slide carries more than \deckcitefootmax of them. A crushed slide
%  should never be a surprise you only discover in the PDF.
\newcount\deck@citefootn
\def\deck@citefootwhere{}
\newcommand{\deck@citefootcount}{%
  \edef\deck@citefoothere{\the\c@page.\the\beamer@slideinframe}%
  \ifx\deck@citefoothere\deck@citefootwhere
    \global\advance\deck@citefootn\@ne
  \else
    \global\deck@citefootn\@ne
    \global\let\deck@citefootwhere\deck@citefoothere
  \fi
  \ifnum\deck@citefootn>\deck@citefootmax\relax
    \PackageWarning{deck-bib}{%
      \the\deck@citefootn\space bottom references on slide
      \the\c@page\space -- more than \the\deck@citefootmax.\MessageBreak
      They will crowd the slide. Use \string\citequiet\space for the
      extras,\MessageBreak
      or \string\deckcitefootoff\space at the top of that frame%
    }%
  \fi
}

% --- Dedup, per slide ------------------------------------------------------
%  Cite the same work twice on one slide and it should appear once at the
%  bottom; cite it again on a later slide and it should appear there too.
%
%  The marker is keyed on the physical page AND on which overlay of the frame
%  is being built. beamer typesets a frame body once per overlay, so a flag
%  reset by a begin-frame hook would fire once and then suppress the block on
%  overlays 2..n -- the reference would vanish the moment you pressed the
%  clicker. Keying on the slide being built needs no hook at all.
\newcommand{\deck@citefootone}[1]{%
  \ifcsname deck@seen@\the\c@page @\the\beamer@slideinframe @#1\endcsname
  \else
    \expandafter\gdef
      \csname deck@seen@\the\c@page @\the\beamer@slideinframe @#1\endcsname{}%
    \deck@citefootcount
    \deck@citefootline{#1}%
  \fi
}

%  \cite{a,b} is legal, so split the list. \forcsvlist trims stray spaces.
\newcommand{\deck@citefoot}[1]{%
  \ifdeck@citefoot
    \forcsvlist{\deck@citefootone}{#1}%
  \fi
}

% --- \cite does both -------------------------------------------------------
%  Wrapped, not reimplemented, so biblatex keeps full control of what [1]
%  actually looks like.
%
%  Deferred to \AtBeginDocument because hyperref -- which beamer loads for us
%  -- patches \cite at the end of the preamble. Grabbing it any earlier would
%  wrap a version that is about to be replaced.
%  biblatex reads ONE optional argument as the postnote and TWO as prenote +
%  postnote, so the wrapper has to hand back exactly as many as it was given.
%  Empty defaults are used rather than -NoValue- markers: an empty argument is
%  something \detokenize can test for certainty, and there is nothing to pass
%  through a second command that might mistake the marker for real text.
\newcommand{\deck@ifempty}[1]{%
  \if\relax\detokenize{#1}\relax
    \expandafter\@firstoftwo
  \else
    \expandafter\@secondoftwo
  \fi
}
\newcommand{\deck@passcite}[3]{%
  \deck@ifempty{#1}%
    {\deck@ifempty{#2}{\deck@origcite{#3}}{\deck@origcite[][#2]{#3}}}%
    {\deck@ifempty{#2}{\deck@origcite[#1]{#3}}{\deck@origcite[#1][#2]{#3}}}%
}

\NewDocumentCommand{\citequiet}{O{} O{} m}{\deck@passcite{#1}{#2}{#3}}

\AtBeginDocument{%
  \let\deck@origcite\cite
  \RenewDocumentCommand{\cite}{s O{} O{} m}{%
    \deck@passcite{#2}{#3}{#4}%
    \IfBooleanF{#1}{\deck@citefoot{#4}}%
  }%
}

\makeatother

%      \addbibresource{refs.bib}
%      ...
%      Shown by \cite{knuth1984}.   ->  "Shown by [1]."  +  [1] at slide bottom
%      \citequiet{knuth1984}        ->  [1] only, nothing at the bottom
%      \cite*{knuth1984}            ->  same as \citequiet, shorter to type
%      \bibliographyframe           ->  the references slide at the end
%
%  \cite takes biblatex's usual optional arguments:
%      \cite[p.~4]{key}          \cite[see][p.~4]{key}
%
%  ---------------------------------------------------------------------------
%  KNOBS -- put any of these in main.tex AFTER % ============================================================================
%  preamble/bib.tex  --  citations and reference list
% ----------------------------------------------------------------------------
%  Numeric citations. \cite{key} does TWO things:
%
%      1. prints  [1]  where you wrote it
%      2. prints the FULL reference, small, at the bottom of THAT slide
%
%  and the same work appears again, with the same number, on the references
%  slide at the end. One command, nothing to keep in sync by hand.
%
%  In your deck:
%      \input{bib}
%      \addbibresource{refs.bib}
%      ...
%      Shown by \cite{knuth1984}.   ->  "Shown by [1]."  +  [1] at slide bottom
%      \citequiet{knuth1984}        ->  [1] only, nothing at the bottom
%      \cite*{knuth1984}            ->  same as \citequiet, shorter to type
%      \bibliographyframe           ->  the references slide at the end
%
%  \cite takes biblatex's usual optional arguments:
%      \cite[p.~4]{key}          \cite[see][p.~4]{key}
%
%  ---------------------------------------------------------------------------
%  KNOBS -- put any of these in main.tex AFTER \input{bib}
%  ---------------------------------------------------------------------------
%      \deckcitefootsize{\scriptsize}  larger bottom text   (default \tiny)
%      \deckcitefootrule{0pt}          drop the hairline above the block
%      \deckcitefootmax{6}             raise the crowding warning threshold
%      \deckcitefootoff                stop adding bottom references
%      \deckcitefooton                 turn them back on
%
%  \deckcitefootoff and \deckcitefooton also work INSIDE a frame, where they
%  apply to that frame only -- useful on a related-work slide with six cites.
%
%  ---------------------------------------------------------------------------
%  WHERE NOT TO PUT \cite
%  ---------------------------------------------------------------------------
%  Not in \frametitle, \section or \caption. Those are moving arguments: the
%  text is written out to the .toc/.aux and typeset again elsewhere, and a
%  footnote cannot survive that. Use \citequiet there.
% ============================================================================

\usepackage[
  backend     = biber,
  style       = numeric,   % [1] in the text, [1] in the reference list
  sorting     = none,      % list ordered by first citation, so [1] really is
                           % the first work you cite -- numbers read in order
  maxbibnames = 6,
  doi         = false,     % slides rarely need these; they eat space
  isbn        = false,
  url         = false,
  eprint      = false,
]{biblatex}

\usepackage{etoolbox}      % \forcsvlist -- biblatex loads it too, this line
                           % is here so the file does not rely on that

% ---------------------------------------------------------------------------
%  THE REFERENCES SLIDE
% ---------------------------------------------------------------------------
%  beamer replaces biblatex's entry label with its own "bibliography item"
%  template. Left empty it drops beamer's generic article icon -- but it also
%  silently deletes the [1] that numeric style depends on. \insertbiblabel
%  puts the number back.
\setbeamertemplate{bibliography item}{\insertbiblabel}
\renewcommand*{\bibfont}{\scriptsize}

% Keep entries from being hyphenated into a mess in a narrow column
\appto{\bibsetup}{\raggedright}

%  A frame that breaks across pages is titled "References I", "References II"
%  and so on. Numbering a list that turned out to fit on one slide just looks
%  like a mistake, so the numeral starts from the second page.
\setbeamertemplate{frametitle continuation}[from second]

% ---------------------------------------------------------------------------
%  A references slide that splits across pages automatically if it is long.
%  Use \bibliographyframe in your deck instead of writing this out by hand.
% ---------------------------------------------------------------------------
\newcommand{\bibliographyframe}[1][References]{%
  \begin{frame}[allowframebreaks]{#1}
    \printbibliography[heading=none]
  \end{frame}%
}

% ===========================================================================
%  THE PER-SLIDE REFERENCE BLOCK
% ===========================================================================
\makeatletter

% --- Knobs -----------------------------------------------------------------
\newcommand*{\deck@citefootsize}{\tiny}
\newcommand*{\deckcitefootsize}[1]{\renewcommand*{\deck@citefootsize}{#1}}

\newlength{\deck@citefootrule}
\setlength{\deck@citefootrule}{0.4pt}
\newcommand*{\deckcitefootrule}[1]{\setlength{\deck@citefootrule}{#1}}

\newcount\deck@citefootmax  \deck@citefootmax=4
\newcommand*{\deckcitefootmax}[1]{\deck@citefootmax=#1\relax}

\newif\ifdeck@citefoot  \deck@citefoottrue
\newcommand*{\deckcitefooton}{\deck@citefoottrue}
\newcommand*{\deckcitefootoff}{\deck@citefootfalse}

% --- The number, on its own ------------------------------------------------
%  Prints just [1] for a key, so the bottom line can be labelled with the
%  same number the reader just saw in the body text.
\DeclareCiteCommand{\deckcitenum}[\mkbibbrackets]
  {}
  {\printfield{labelnumber}}
  {\multicitedelim}
  {}

% --- The hairline above the block ------------------------------------------
%  This is LaTeX's footnote rule, so an ordinary \footnote gets the same line.
%  \deckcitefootrule{0pt} makes it disappear without disturbing the spacing.
\renewcommand{\footnoterule}{%
  \kern-3pt%
  \hbox{\color{BrandRule}%
        \vrule width 0.35\textwidth height \deck@citefootrule depth 0pt}%
  \kern2.6pt%
}

% --- One bottom line -------------------------------------------------------
%  An unmarked footnote: beamer already places footnotes at the bottom of the
%  frame, above the footer, and keeps them working under [allowframebreaks].
%  The mark is emptied and the footnote counter rolled back, because the line
%  carries its own [n] instead.
%
%  The footnote template is overridden LOCALLY. beamer's default reserves
%  1.8em for the mark; with no mark that is 1.8em of dead indent on every
%  line. Ordinary footnotes elsewhere in the deck are untouched.
%
%  The four \setbeamerfont lines are not optional. \fullcite goes through
%  beamer's bibliography drivers, which apply "bibliography entry ..." fonts --
%  absolute sizes that would override \deckcitefootsize and leave the [n]
%  label tiny next to a much larger reference. Unstarred \setbeamerfont
%  changes only the size and leaves series and shape alone, and all four are
%  local to the group, so the references slide at the end is unaffected.
\newcommand{\deck@citefootline}[1]{%
  \begingroup
    \setbeamerfont{bibliography entry author}  {size=\deck@citefootsize}%
    \setbeamerfont{bibliography entry title}   {size=\deck@citefootsize}%
    \setbeamerfont{bibliography entry location}{size=\deck@citefootsize}%
    \setbeamerfont{bibliography entry note}    {size=\deck@citefootsize}%
    \setbeamertemplate{footnote}{%
      \parindent0pt\noindent\raggedright
      \insertfootnotetext\par
    }%
    \renewcommand{\thefootnote}{}%
    \footnote{{\deck@citefootsize\deckcitenum{#1}\ \fullcite{#1}}}%
    \addtocounter{footnote}{-1}%
  \endgroup
}

% --- Crowding warning ------------------------------------------------------
%  Counts bottom references on the current slide and says so in the log once
%  the slide carries more than \deckcitefootmax of them. A crushed slide
%  should never be a surprise you only discover in the PDF.
\newcount\deck@citefootn
\def\deck@citefootwhere{}
\newcommand{\deck@citefootcount}{%
  \edef\deck@citefoothere{\the\c@page.\the\beamer@slideinframe}%
  \ifx\deck@citefoothere\deck@citefootwhere
    \global\advance\deck@citefootn\@ne
  \else
    \global\deck@citefootn\@ne
    \global\let\deck@citefootwhere\deck@citefoothere
  \fi
  \ifnum\deck@citefootn>\deck@citefootmax\relax
    \PackageWarning{deck-bib}{%
      \the\deck@citefootn\space bottom references on slide
      \the\c@page\space -- more than \the\deck@citefootmax.\MessageBreak
      They will crowd the slide. Use \string\citequiet\space for the
      extras,\MessageBreak
      or \string\deckcitefootoff\space at the top of that frame%
    }%
  \fi
}

% --- Dedup, per slide ------------------------------------------------------
%  Cite the same work twice on one slide and it should appear once at the
%  bottom; cite it again on a later slide and it should appear there too.
%
%  The marker is keyed on the physical page AND on which overlay of the frame
%  is being built. beamer typesets a frame body once per overlay, so a flag
%  reset by a begin-frame hook would fire once and then suppress the block on
%  overlays 2..n -- the reference would vanish the moment you pressed the
%  clicker. Keying on the slide being built needs no hook at all.
\newcommand{\deck@citefootone}[1]{%
  \ifcsname deck@seen@\the\c@page @\the\beamer@slideinframe @#1\endcsname
  \else
    \expandafter\gdef
      \csname deck@seen@\the\c@page @\the\beamer@slideinframe @#1\endcsname{}%
    \deck@citefootcount
    \deck@citefootline{#1}%
  \fi
}

%  \cite{a,b} is legal, so split the list. \forcsvlist trims stray spaces.
\newcommand{\deck@citefoot}[1]{%
  \ifdeck@citefoot
    \forcsvlist{\deck@citefootone}{#1}%
  \fi
}

% --- \cite does both -------------------------------------------------------
%  Wrapped, not reimplemented, so biblatex keeps full control of what [1]
%  actually looks like.
%
%  Deferred to \AtBeginDocument because hyperref -- which beamer loads for us
%  -- patches \cite at the end of the preamble. Grabbing it any earlier would
%  wrap a version that is about to be replaced.
%  biblatex reads ONE optional argument as the postnote and TWO as prenote +
%  postnote, so the wrapper has to hand back exactly as many as it was given.
%  Empty defaults are used rather than -NoValue- markers: an empty argument is
%  something \detokenize can test for certainty, and there is nothing to pass
%  through a second command that might mistake the marker for real text.
\newcommand{\deck@ifempty}[1]{%
  \if\relax\detokenize{#1}\relax
    \expandafter\@firstoftwo
  \else
    \expandafter\@secondoftwo
  \fi
}
\newcommand{\deck@passcite}[3]{%
  \deck@ifempty{#1}%
    {\deck@ifempty{#2}{\deck@origcite{#3}}{\deck@origcite[][#2]{#3}}}%
    {\deck@ifempty{#2}{\deck@origcite[#1]{#3}}{\deck@origcite[#1][#2]{#3}}}%
}

\NewDocumentCommand{\citequiet}{O{} O{} m}{\deck@passcite{#1}{#2}{#3}}

\AtBeginDocument{%
  \let\deck@origcite\cite
  \RenewDocumentCommand{\cite}{s O{} O{} m}{%
    \deck@passcite{#2}{#3}{#4}%
    \IfBooleanF{#1}{\deck@citefoot{#4}}%
  }%
}

\makeatother

%  ---------------------------------------------------------------------------
%      \deckcitefootsize{\scriptsize}  larger bottom text   (default \tiny)
%      \deckcitefootrule{0pt}          drop the hairline above the block
%      \deckcitefootmax{6}             raise the crowding warning threshold
%      \deckcitefootoff                stop adding bottom references
%      \deckcitefooton                 turn them back on
%
%  \deckcitefootoff and \deckcitefooton also work INSIDE a frame, where they
%  apply to that frame only -- useful on a related-work slide with six cites.
%
%  ---------------------------------------------------------------------------
%  WHERE NOT TO PUT \cite
%  ---------------------------------------------------------------------------
%  Not in \frametitle, \section or \caption. Those are moving arguments: the
%  text is written out to the .toc/.aux and typeset again elsewhere, and a
%  footnote cannot survive that. Use \citequiet there.
% ============================================================================

\usepackage[
  backend     = biber,
  style       = numeric,   % [1] in the text, [1] in the reference list
  sorting     = none,      % list ordered by first citation, so [1] really is
                           % the first work you cite -- numbers read in order
  maxbibnames = 6,
  doi         = false,     % slides rarely need these; they eat space
  isbn        = false,
  url         = false,
  eprint      = false,
]{biblatex}

\usepackage{etoolbox}      % \forcsvlist -- biblatex loads it too, this line
                           % is here so the file does not rely on that

% ---------------------------------------------------------------------------
%  THE REFERENCES SLIDE
% ---------------------------------------------------------------------------
%  beamer replaces biblatex's entry label with its own "bibliography item"
%  template. Left empty it drops beamer's generic article icon -- but it also
%  silently deletes the [1] that numeric style depends on. \insertbiblabel
%  puts the number back.
\setbeamertemplate{bibliography item}{\insertbiblabel}
\renewcommand*{\bibfont}{\scriptsize}

% Keep entries from being hyphenated into a mess in a narrow column
\appto{\bibsetup}{\raggedright}

%  A frame that breaks across pages is titled "References I", "References II"
%  and so on. Numbering a list that turned out to fit on one slide just looks
%  like a mistake, so the numeral starts from the second page.
\setbeamertemplate{frametitle continuation}[from second]

% ---------------------------------------------------------------------------
%  A references slide that splits across pages automatically if it is long.
%  Use \bibliographyframe in your deck instead of writing this out by hand.
% ---------------------------------------------------------------------------
\newcommand{\bibliographyframe}[1][References]{%
  \begin{frame}[allowframebreaks]{#1}
    \printbibliography[heading=none]
  \end{frame}%
}

% ===========================================================================
%  THE PER-SLIDE REFERENCE BLOCK
% ===========================================================================
\makeatletter

% --- Knobs -----------------------------------------------------------------
\newcommand*{\deck@citefootsize}{\tiny}
\newcommand*{\deckcitefootsize}[1]{\renewcommand*{\deck@citefootsize}{#1}}

\newlength{\deck@citefootrule}
\setlength{\deck@citefootrule}{0.4pt}
\newcommand*{\deckcitefootrule}[1]{\setlength{\deck@citefootrule}{#1}}

\newcount\deck@citefootmax  \deck@citefootmax=4
\newcommand*{\deckcitefootmax}[1]{\deck@citefootmax=#1\relax}

\newif\ifdeck@citefoot  \deck@citefoottrue
\newcommand*{\deckcitefooton}{\deck@citefoottrue}
\newcommand*{\deckcitefootoff}{\deck@citefootfalse}

% --- The number, on its own ------------------------------------------------
%  Prints just [1] for a key, so the bottom line can be labelled with the
%  same number the reader just saw in the body text.
\DeclareCiteCommand{\deckcitenum}[\mkbibbrackets]
  {}
  {\printfield{labelnumber}}
  {\multicitedelim}
  {}

% --- The hairline above the block ------------------------------------------
%  This is LaTeX's footnote rule, so an ordinary \footnote gets the same line.
%  \deckcitefootrule{0pt} makes it disappear without disturbing the spacing.
\renewcommand{\footnoterule}{%
  \kern-3pt%
  \hbox{\color{BrandRule}%
        \vrule width 0.35\textwidth height \deck@citefootrule depth 0pt}%
  \kern2.6pt%
}

% --- One bottom line -------------------------------------------------------
%  An unmarked footnote: beamer already places footnotes at the bottom of the
%  frame, above the footer, and keeps them working under [allowframebreaks].
%  The mark is emptied and the footnote counter rolled back, because the line
%  carries its own [n] instead.
%
%  The footnote template is overridden LOCALLY. beamer's default reserves
%  1.8em for the mark; with no mark that is 1.8em of dead indent on every
%  line. Ordinary footnotes elsewhere in the deck are untouched.
%
%  The four \setbeamerfont lines are not optional. \fullcite goes through
%  beamer's bibliography drivers, which apply "bibliography entry ..." fonts --
%  absolute sizes that would override \deckcitefootsize and leave the [n]
%  label tiny next to a much larger reference. Unstarred \setbeamerfont
%  changes only the size and leaves series and shape alone, and all four are
%  local to the group, so the references slide at the end is unaffected.
\newcommand{\deck@citefootline}[1]{%
  \begingroup
    \setbeamerfont{bibliography entry author}  {size=\deck@citefootsize}%
    \setbeamerfont{bibliography entry title}   {size=\deck@citefootsize}%
    \setbeamerfont{bibliography entry location}{size=\deck@citefootsize}%
    \setbeamerfont{bibliography entry note}    {size=\deck@citefootsize}%
    \setbeamertemplate{footnote}{%
      \parindent0pt\noindent\raggedright
      \insertfootnotetext\par
    }%
    \renewcommand{\thefootnote}{}%
    \footnote{{\deck@citefootsize\deckcitenum{#1}\ \fullcite{#1}}}%
    \addtocounter{footnote}{-1}%
  \endgroup
}

% --- Crowding warning ------------------------------------------------------
%  Counts bottom references on the current slide and says so in the log once
%  the slide carries more than \deckcitefootmax of them. A crushed slide
%  should never be a surprise you only discover in the PDF.
\newcount\deck@citefootn
\def\deck@citefootwhere{}
\newcommand{\deck@citefootcount}{%
  \edef\deck@citefoothere{\the\c@page.\the\beamer@slideinframe}%
  \ifx\deck@citefoothere\deck@citefootwhere
    \global\advance\deck@citefootn\@ne
  \else
    \global\deck@citefootn\@ne
    \global\let\deck@citefootwhere\deck@citefoothere
  \fi
  \ifnum\deck@citefootn>\deck@citefootmax\relax
    \PackageWarning{deck-bib}{%
      \the\deck@citefootn\space bottom references on slide
      \the\c@page\space -- more than \the\deck@citefootmax.\MessageBreak
      They will crowd the slide. Use \string\citequiet\space for the
      extras,\MessageBreak
      or \string\deckcitefootoff\space at the top of that frame%
    }%
  \fi
}

% --- Dedup, per slide ------------------------------------------------------
%  Cite the same work twice on one slide and it should appear once at the
%  bottom; cite it again on a later slide and it should appear there too.
%
%  The marker is keyed on the physical page AND on which overlay of the frame
%  is being built. beamer typesets a frame body once per overlay, so a flag
%  reset by a begin-frame hook would fire once and then suppress the block on
%  overlays 2..n -- the reference would vanish the moment you pressed the
%  clicker. Keying on the slide being built needs no hook at all.
\newcommand{\deck@citefootone}[1]{%
  \ifcsname deck@seen@\the\c@page @\the\beamer@slideinframe @#1\endcsname
  \else
    \expandafter\gdef
      \csname deck@seen@\the\c@page @\the\beamer@slideinframe @#1\endcsname{}%
    \deck@citefootcount
    \deck@citefootline{#1}%
  \fi
}

%  \cite{a,b} is legal, so split the list. \forcsvlist trims stray spaces.
\newcommand{\deck@citefoot}[1]{%
  \ifdeck@citefoot
    \forcsvlist{\deck@citefootone}{#1}%
  \fi
}

% --- \cite does both -------------------------------------------------------
%  Wrapped, not reimplemented, so biblatex keeps full control of what [1]
%  actually looks like.
%
%  Deferred to \AtBeginDocument because hyperref -- which beamer loads for us
%  -- patches \cite at the end of the preamble. Grabbing it any earlier would
%  wrap a version that is about to be replaced.
%  biblatex reads ONE optional argument as the postnote and TWO as prenote +
%  postnote, so the wrapper has to hand back exactly as many as it was given.
%  Empty defaults are used rather than -NoValue- markers: an empty argument is
%  something \detokenize can test for certainty, and there is nothing to pass
%  through a second command that might mistake the marker for real text.
\newcommand{\deck@ifempty}[1]{%
  \if\relax\detokenize{#1}\relax
    \expandafter\@firstoftwo
  \else
    \expandafter\@secondoftwo
  \fi
}
\newcommand{\deck@passcite}[3]{%
  \deck@ifempty{#1}%
    {\deck@ifempty{#2}{\deck@origcite{#3}}{\deck@origcite[][#2]{#3}}}%
    {\deck@ifempty{#2}{\deck@origcite[#1]{#3}}{\deck@origcite[#1][#2]{#3}}}%
}

\NewDocumentCommand{\citequiet}{O{} O{} m}{\deck@passcite{#1}{#2}{#3}}

\AtBeginDocument{%
  \let\deck@origcite\cite
  \RenewDocumentCommand{\cite}{s O{} O{} m}{%
    \deck@passcite{#2}{#3}{#4}%
    \IfBooleanF{#1}{\deck@citefoot{#4}}%
  }%
}

\makeatother

%  ---------------------------------------------------------------------------
%      \deckcitefootsize{\scriptsize}  larger bottom text   (default \tiny)
%      \deckcitefootrule{0pt}          drop the hairline above the block
%      \deckcitefootmax{6}             raise the crowding warning threshold
%      \deckcitefootoff                stop adding bottom references
%      \deckcitefooton                 turn them back on
%
%  \deckcitefootoff and \deckcitefooton also work INSIDE a frame, where they
%  apply to that frame only -- useful on a related-work slide with six cites.
%
%  ---------------------------------------------------------------------------
%  WHERE NOT TO PUT \cite
%  ---------------------------------------------------------------------------
%  Not in \frametitle, \section or \caption. Those are moving arguments: the
%  text is written out to the .toc/.aux and typeset again elsewhere, and a
%  footnote cannot survive that. Use \citequiet there.
% ============================================================================

\usepackage[
  backend     = biber,
  style       = numeric,   % [1] in the text, [1] in the reference list
  sorting     = none,      % list ordered by first citation, so [1] really is
                           % the first work you cite -- numbers read in order
  maxbibnames = 6,
  doi         = false,     % slides rarely need these; they eat space
  isbn        = false,
  url         = false,
  eprint      = false,
]{biblatex}

\usepackage{etoolbox}      % \forcsvlist -- biblatex loads it too, this line
                           % is here so the file does not rely on that

% ---------------------------------------------------------------------------
%  THE REFERENCES SLIDE
% ---------------------------------------------------------------------------
%  beamer replaces biblatex's entry label with its own "bibliography item"
%  template. Left empty it drops beamer's generic article icon -- but it also
%  silently deletes the [1] that numeric style depends on. \insertbiblabel
%  puts the number back.
\setbeamertemplate{bibliography item}{\insertbiblabel}
\renewcommand*{\bibfont}{\scriptsize}

% Keep entries from being hyphenated into a mess in a narrow column
\appto{\bibsetup}{\raggedright}

%  A frame that breaks across pages is titled "References I", "References II"
%  and so on. Numbering a list that turned out to fit on one slide just looks
%  like a mistake, so the numeral starts from the second page.
\setbeamertemplate{frametitle continuation}[from second]

% ---------------------------------------------------------------------------
%  A references slide that splits across pages automatically if it is long.
%  Use \bibliographyframe in your deck instead of writing this out by hand.
% ---------------------------------------------------------------------------
\newcommand{\bibliographyframe}[1][References]{%
  \begin{frame}[allowframebreaks]{#1}
    \printbibliography[heading=none]
  \end{frame}%
}

% ===========================================================================
%  THE PER-SLIDE REFERENCE BLOCK
% ===========================================================================
\makeatletter

% --- Knobs -----------------------------------------------------------------
\newcommand*{\deck@citefootsize}{\tiny}
\newcommand*{\deckcitefootsize}[1]{\renewcommand*{\deck@citefootsize}{#1}}

\newlength{\deck@citefootrule}
\setlength{\deck@citefootrule}{0.4pt}
\newcommand*{\deckcitefootrule}[1]{\setlength{\deck@citefootrule}{#1}}

\newcount\deck@citefootmax  \deck@citefootmax=4
\newcommand*{\deckcitefootmax}[1]{\deck@citefootmax=#1\relax}

\newif\ifdeck@citefoot  \deck@citefoottrue
\newcommand*{\deckcitefooton}{\deck@citefoottrue}
\newcommand*{\deckcitefootoff}{\deck@citefootfalse}

% --- The number, on its own ------------------------------------------------
%  Prints just [1] for a key, so the bottom line can be labelled with the
%  same number the reader just saw in the body text.
\DeclareCiteCommand{\deckcitenum}[\mkbibbrackets]
  {}
  {\printfield{labelnumber}}
  {\multicitedelim}
  {}

% --- The hairline above the block ------------------------------------------
%  This is LaTeX's footnote rule, so an ordinary \footnote gets the same line.
%  \deckcitefootrule{0pt} makes it disappear without disturbing the spacing.
\renewcommand{\footnoterule}{%
  \kern-3pt%
  \hbox{\color{BrandRule}%
        \vrule width 0.35\textwidth height \deck@citefootrule depth 0pt}%
  \kern2.6pt%
}

% --- One bottom line -------------------------------------------------------
%  An unmarked footnote: beamer already places footnotes at the bottom of the
%  frame, above the footer, and keeps them working under [allowframebreaks].
%  The mark is emptied and the footnote counter rolled back, because the line
%  carries its own [n] instead.
%
%  The footnote template is overridden LOCALLY. beamer's default reserves
%  1.8em for the mark; with no mark that is 1.8em of dead indent on every
%  line. Ordinary footnotes elsewhere in the deck are untouched.
%
%  The four \setbeamerfont lines are not optional. \fullcite goes through
%  beamer's bibliography drivers, which apply "bibliography entry ..." fonts --
%  absolute sizes that would override \deckcitefootsize and leave the [n]
%  label tiny next to a much larger reference. Unstarred \setbeamerfont
%  changes only the size and leaves series and shape alone, and all four are
%  local to the group, so the references slide at the end is unaffected.
\newcommand{\deck@citefootline}[1]{%
  \begingroup
    \setbeamerfont{bibliography entry author}  {size=\deck@citefootsize}%
    \setbeamerfont{bibliography entry title}   {size=\deck@citefootsize}%
    \setbeamerfont{bibliography entry location}{size=\deck@citefootsize}%
    \setbeamerfont{bibliography entry note}    {size=\deck@citefootsize}%
    \setbeamertemplate{footnote}{%
      \parindent0pt\noindent\raggedright
      \insertfootnotetext\par
    }%
    \renewcommand{\thefootnote}{}%
    \footnote{{\deck@citefootsize\deckcitenum{#1}\ \fullcite{#1}}}%
    \addtocounter{footnote}{-1}%
  \endgroup
}

% --- Crowding warning ------------------------------------------------------
%  Counts bottom references on the current slide and says so in the log once
%  the slide carries more than \deckcitefootmax of them. A crushed slide
%  should never be a surprise you only discover in the PDF.
\newcount\deck@citefootn
\def\deck@citefootwhere{}
\newcommand{\deck@citefootcount}{%
  \edef\deck@citefoothere{\the\c@page.\the\beamer@slideinframe}%
  \ifx\deck@citefoothere\deck@citefootwhere
    \global\advance\deck@citefootn\@ne
  \else
    \global\deck@citefootn\@ne
    \global\let\deck@citefootwhere\deck@citefoothere
  \fi
  \ifnum\deck@citefootn>\deck@citefootmax\relax
    \PackageWarning{deck-bib}{%
      \the\deck@citefootn\space bottom references on slide
      \the\c@page\space -- more than \the\deck@citefootmax.\MessageBreak
      They will crowd the slide. Use \string\citequiet\space for the
      extras,\MessageBreak
      or \string\deckcitefootoff\space at the top of that frame%
    }%
  \fi
}

% --- Dedup, per slide ------------------------------------------------------
%  Cite the same work twice on one slide and it should appear once at the
%  bottom; cite it again on a later slide and it should appear there too.
%
%  The marker is keyed on the physical page AND on which overlay of the frame
%  is being built. beamer typesets a frame body once per overlay, so a flag
%  reset by a begin-frame hook would fire once and then suppress the block on
%  overlays 2..n -- the reference would vanish the moment you pressed the
%  clicker. Keying on the slide being built needs no hook at all.
\newcommand{\deck@citefootone}[1]{%
  \ifcsname deck@seen@\the\c@page @\the\beamer@slideinframe @#1\endcsname
  \else
    \expandafter\gdef
      \csname deck@seen@\the\c@page @\the\beamer@slideinframe @#1\endcsname{}%
    \deck@citefootcount
    \deck@citefootline{#1}%
  \fi
}

%  \cite{a,b} is legal, so split the list. \forcsvlist trims stray spaces.
\newcommand{\deck@citefoot}[1]{%
  \ifdeck@citefoot
    \forcsvlist{\deck@citefootone}{#1}%
  \fi
}

% --- \cite does both -------------------------------------------------------
%  Wrapped, not reimplemented, so biblatex keeps full control of what [1]
%  actually looks like.
%
%  Deferred to \AtBeginDocument because hyperref -- which beamer loads for us
%  -- patches \cite at the end of the preamble. Grabbing it any earlier would
%  wrap a version that is about to be replaced.
%  biblatex reads ONE optional argument as the postnote and TWO as prenote +
%  postnote, so the wrapper has to hand back exactly as many as it was given.
%  Empty defaults are used rather than -NoValue- markers: an empty argument is
%  something \detokenize can test for certainty, and there is nothing to pass
%  through a second command that might mistake the marker for real text.
\newcommand{\deck@ifempty}[1]{%
  \if\relax\detokenize{#1}\relax
    \expandafter\@firstoftwo
  \else
    \expandafter\@secondoftwo
  \fi
}
\newcommand{\deck@passcite}[3]{%
  \deck@ifempty{#1}%
    {\deck@ifempty{#2}{\deck@origcite{#3}}{\deck@origcite[][#2]{#3}}}%
    {\deck@ifempty{#2}{\deck@origcite[#1]{#3}}{\deck@origcite[#1][#2]{#3}}}%
}

\NewDocumentCommand{\citequiet}{O{} O{} m}{\deck@passcite{#1}{#2}{#3}}

\AtBeginDocument{%
  \let\deck@origcite\cite
  \RenewDocumentCommand{\cite}{s O{} O{} m}{%
    \deck@passcite{#2}{#3}{#4}%
    \IfBooleanF{#1}{\deck@citefoot{#4}}%
  }%
}

\makeatother

%      \addbibresource{refs.bib}
%      ...
%      Shown by \cite{knuth1984}.   ->  "Shown by [1]."  +  [1] at slide bottom
%      \citequiet{knuth1984}        ->  [1] only, nothing at the bottom
%      \cite*{knuth1984}            ->  same as \citequiet, shorter to type
%      \bibliographyframe           ->  the references slide at the end
%
%  \cite takes biblatex's usual optional arguments:
%      \cite[p.~4]{key}          \cite[see][p.~4]{key}
%
%  ---------------------------------------------------------------------------
%  KNOBS -- put any of these in main.tex AFTER % ============================================================================
%  preamble/bib.tex  --  citations and reference list
% ----------------------------------------------------------------------------
%  Numeric citations. \cite{key} does TWO things:
%
%      1. prints  [1]  where you wrote it
%      2. prints the FULL reference, small, at the bottom of THAT slide
%
%  and the same work appears again, with the same number, on the references
%  slide at the end. One command, nothing to keep in sync by hand.
%
%  In your deck:
%      % ============================================================================
%  preamble/bib.tex  --  citations and reference list
% ----------------------------------------------------------------------------
%  Numeric citations. \cite{key} does TWO things:
%
%      1. prints  [1]  where you wrote it
%      2. prints the FULL reference, small, at the bottom of THAT slide
%
%  and the same work appears again, with the same number, on the references
%  slide at the end. One command, nothing to keep in sync by hand.
%
%  In your deck:
%      % ============================================================================
%  preamble/bib.tex  --  citations and reference list
% ----------------------------------------------------------------------------
%  Numeric citations. \cite{key} does TWO things:
%
%      1. prints  [1]  where you wrote it
%      2. prints the FULL reference, small, at the bottom of THAT slide
%
%  and the same work appears again, with the same number, on the references
%  slide at the end. One command, nothing to keep in sync by hand.
%
%  In your deck:
%      \input{bib}
%      \addbibresource{refs.bib}
%      ...
%      Shown by \cite{knuth1984}.   ->  "Shown by [1]."  +  [1] at slide bottom
%      \citequiet{knuth1984}        ->  [1] only, nothing at the bottom
%      \cite*{knuth1984}            ->  same as \citequiet, shorter to type
%      \bibliographyframe           ->  the references slide at the end
%
%  \cite takes biblatex's usual optional arguments:
%      \cite[p.~4]{key}          \cite[see][p.~4]{key}
%
%  ---------------------------------------------------------------------------
%  KNOBS -- put any of these in main.tex AFTER \input{bib}
%  ---------------------------------------------------------------------------
%      \deckcitefootsize{\scriptsize}  larger bottom text   (default \tiny)
%      \deckcitefootrule{0pt}          drop the hairline above the block
%      \deckcitefootmax{6}             raise the crowding warning threshold
%      \deckcitefootoff                stop adding bottom references
%      \deckcitefooton                 turn them back on
%
%  \deckcitefootoff and \deckcitefooton also work INSIDE a frame, where they
%  apply to that frame only -- useful on a related-work slide with six cites.
%
%  ---------------------------------------------------------------------------
%  WHERE NOT TO PUT \cite
%  ---------------------------------------------------------------------------
%  Not in \frametitle, \section or \caption. Those are moving arguments: the
%  text is written out to the .toc/.aux and typeset again elsewhere, and a
%  footnote cannot survive that. Use \citequiet there.
% ============================================================================

\usepackage[
  backend     = biber,
  style       = numeric,   % [1] in the text, [1] in the reference list
  sorting     = none,      % list ordered by first citation, so [1] really is
                           % the first work you cite -- numbers read in order
  maxbibnames = 6,
  doi         = false,     % slides rarely need these; they eat space
  isbn        = false,
  url         = false,
  eprint      = false,
]{biblatex}

\usepackage{etoolbox}      % \forcsvlist -- biblatex loads it too, this line
                           % is here so the file does not rely on that

% ---------------------------------------------------------------------------
%  THE REFERENCES SLIDE
% ---------------------------------------------------------------------------
%  beamer replaces biblatex's entry label with its own "bibliography item"
%  template. Left empty it drops beamer's generic article icon -- but it also
%  silently deletes the [1] that numeric style depends on. \insertbiblabel
%  puts the number back.
\setbeamertemplate{bibliography item}{\insertbiblabel}
\renewcommand*{\bibfont}{\scriptsize}

% Keep entries from being hyphenated into a mess in a narrow column
\appto{\bibsetup}{\raggedright}

%  A frame that breaks across pages is titled "References I", "References II"
%  and so on. Numbering a list that turned out to fit on one slide just looks
%  like a mistake, so the numeral starts from the second page.
\setbeamertemplate{frametitle continuation}[from second]

% ---------------------------------------------------------------------------
%  A references slide that splits across pages automatically if it is long.
%  Use \bibliographyframe in your deck instead of writing this out by hand.
% ---------------------------------------------------------------------------
\newcommand{\bibliographyframe}[1][References]{%
  \begin{frame}[allowframebreaks]{#1}
    \printbibliography[heading=none]
  \end{frame}%
}

% ===========================================================================
%  THE PER-SLIDE REFERENCE BLOCK
% ===========================================================================
\makeatletter

% --- Knobs -----------------------------------------------------------------
\newcommand*{\deck@citefootsize}{\tiny}
\newcommand*{\deckcitefootsize}[1]{\renewcommand*{\deck@citefootsize}{#1}}

\newlength{\deck@citefootrule}
\setlength{\deck@citefootrule}{0.4pt}
\newcommand*{\deckcitefootrule}[1]{\setlength{\deck@citefootrule}{#1}}

\newcount\deck@citefootmax  \deck@citefootmax=4
\newcommand*{\deckcitefootmax}[1]{\deck@citefootmax=#1\relax}

\newif\ifdeck@citefoot  \deck@citefoottrue
\newcommand*{\deckcitefooton}{\deck@citefoottrue}
\newcommand*{\deckcitefootoff}{\deck@citefootfalse}

% --- The number, on its own ------------------------------------------------
%  Prints just [1] for a key, so the bottom line can be labelled with the
%  same number the reader just saw in the body text.
\DeclareCiteCommand{\deckcitenum}[\mkbibbrackets]
  {}
  {\printfield{labelnumber}}
  {\multicitedelim}
  {}

% --- The hairline above the block ------------------------------------------
%  This is LaTeX's footnote rule, so an ordinary \footnote gets the same line.
%  \deckcitefootrule{0pt} makes it disappear without disturbing the spacing.
\renewcommand{\footnoterule}{%
  \kern-3pt%
  \hbox{\color{BrandRule}%
        \vrule width 0.35\textwidth height \deck@citefootrule depth 0pt}%
  \kern2.6pt%
}

% --- One bottom line -------------------------------------------------------
%  An unmarked footnote: beamer already places footnotes at the bottom of the
%  frame, above the footer, and keeps them working under [allowframebreaks].
%  The mark is emptied and the footnote counter rolled back, because the line
%  carries its own [n] instead.
%
%  The footnote template is overridden LOCALLY. beamer's default reserves
%  1.8em for the mark; with no mark that is 1.8em of dead indent on every
%  line. Ordinary footnotes elsewhere in the deck are untouched.
%
%  The four \setbeamerfont lines are not optional. \fullcite goes through
%  beamer's bibliography drivers, which apply "bibliography entry ..." fonts --
%  absolute sizes that would override \deckcitefootsize and leave the [n]
%  label tiny next to a much larger reference. Unstarred \setbeamerfont
%  changes only the size and leaves series and shape alone, and all four are
%  local to the group, so the references slide at the end is unaffected.
\newcommand{\deck@citefootline}[1]{%
  \begingroup
    \setbeamerfont{bibliography entry author}  {size=\deck@citefootsize}%
    \setbeamerfont{bibliography entry title}   {size=\deck@citefootsize}%
    \setbeamerfont{bibliography entry location}{size=\deck@citefootsize}%
    \setbeamerfont{bibliography entry note}    {size=\deck@citefootsize}%
    \setbeamertemplate{footnote}{%
      \parindent0pt\noindent\raggedright
      \insertfootnotetext\par
    }%
    \renewcommand{\thefootnote}{}%
    \footnote{{\deck@citefootsize\deckcitenum{#1}\ \fullcite{#1}}}%
    \addtocounter{footnote}{-1}%
  \endgroup
}

% --- Crowding warning ------------------------------------------------------
%  Counts bottom references on the current slide and says so in the log once
%  the slide carries more than \deckcitefootmax of them. A crushed slide
%  should never be a surprise you only discover in the PDF.
\newcount\deck@citefootn
\def\deck@citefootwhere{}
\newcommand{\deck@citefootcount}{%
  \edef\deck@citefoothere{\the\c@page.\the\beamer@slideinframe}%
  \ifx\deck@citefoothere\deck@citefootwhere
    \global\advance\deck@citefootn\@ne
  \else
    \global\deck@citefootn\@ne
    \global\let\deck@citefootwhere\deck@citefoothere
  \fi
  \ifnum\deck@citefootn>\deck@citefootmax\relax
    \PackageWarning{deck-bib}{%
      \the\deck@citefootn\space bottom references on slide
      \the\c@page\space -- more than \the\deck@citefootmax.\MessageBreak
      They will crowd the slide. Use \string\citequiet\space for the
      extras,\MessageBreak
      or \string\deckcitefootoff\space at the top of that frame%
    }%
  \fi
}

% --- Dedup, per slide ------------------------------------------------------
%  Cite the same work twice on one slide and it should appear once at the
%  bottom; cite it again on a later slide and it should appear there too.
%
%  The marker is keyed on the physical page AND on which overlay of the frame
%  is being built. beamer typesets a frame body once per overlay, so a flag
%  reset by a begin-frame hook would fire once and then suppress the block on
%  overlays 2..n -- the reference would vanish the moment you pressed the
%  clicker. Keying on the slide being built needs no hook at all.
\newcommand{\deck@citefootone}[1]{%
  \ifcsname deck@seen@\the\c@page @\the\beamer@slideinframe @#1\endcsname
  \else
    \expandafter\gdef
      \csname deck@seen@\the\c@page @\the\beamer@slideinframe @#1\endcsname{}%
    \deck@citefootcount
    \deck@citefootline{#1}%
  \fi
}

%  \cite{a,b} is legal, so split the list. \forcsvlist trims stray spaces.
\newcommand{\deck@citefoot}[1]{%
  \ifdeck@citefoot
    \forcsvlist{\deck@citefootone}{#1}%
  \fi
}

% --- \cite does both -------------------------------------------------------
%  Wrapped, not reimplemented, so biblatex keeps full control of what [1]
%  actually looks like.
%
%  Deferred to \AtBeginDocument because hyperref -- which beamer loads for us
%  -- patches \cite at the end of the preamble. Grabbing it any earlier would
%  wrap a version that is about to be replaced.
%  biblatex reads ONE optional argument as the postnote and TWO as prenote +
%  postnote, so the wrapper has to hand back exactly as many as it was given.
%  Empty defaults are used rather than -NoValue- markers: an empty argument is
%  something \detokenize can test for certainty, and there is nothing to pass
%  through a second command that might mistake the marker for real text.
\newcommand{\deck@ifempty}[1]{%
  \if\relax\detokenize{#1}\relax
    \expandafter\@firstoftwo
  \else
    \expandafter\@secondoftwo
  \fi
}
\newcommand{\deck@passcite}[3]{%
  \deck@ifempty{#1}%
    {\deck@ifempty{#2}{\deck@origcite{#3}}{\deck@origcite[][#2]{#3}}}%
    {\deck@ifempty{#2}{\deck@origcite[#1]{#3}}{\deck@origcite[#1][#2]{#3}}}%
}

\NewDocumentCommand{\citequiet}{O{} O{} m}{\deck@passcite{#1}{#2}{#3}}

\AtBeginDocument{%
  \let\deck@origcite\cite
  \RenewDocumentCommand{\cite}{s O{} O{} m}{%
    \deck@passcite{#2}{#3}{#4}%
    \IfBooleanF{#1}{\deck@citefoot{#4}}%
  }%
}

\makeatother

%      \addbibresource{refs.bib}
%      ...
%      Shown by \cite{knuth1984}.   ->  "Shown by [1]."  +  [1] at slide bottom
%      \citequiet{knuth1984}        ->  [1] only, nothing at the bottom
%      \cite*{knuth1984}            ->  same as \citequiet, shorter to type
%      \bibliographyframe           ->  the references slide at the end
%
%  \cite takes biblatex's usual optional arguments:
%      \cite[p.~4]{key}          \cite[see][p.~4]{key}
%
%  ---------------------------------------------------------------------------
%  KNOBS -- put any of these in main.tex AFTER % ============================================================================
%  preamble/bib.tex  --  citations and reference list
% ----------------------------------------------------------------------------
%  Numeric citations. \cite{key} does TWO things:
%
%      1. prints  [1]  where you wrote it
%      2. prints the FULL reference, small, at the bottom of THAT slide
%
%  and the same work appears again, with the same number, on the references
%  slide at the end. One command, nothing to keep in sync by hand.
%
%  In your deck:
%      \input{bib}
%      \addbibresource{refs.bib}
%      ...
%      Shown by \cite{knuth1984}.   ->  "Shown by [1]."  +  [1] at slide bottom
%      \citequiet{knuth1984}        ->  [1] only, nothing at the bottom
%      \cite*{knuth1984}            ->  same as \citequiet, shorter to type
%      \bibliographyframe           ->  the references slide at the end
%
%  \cite takes biblatex's usual optional arguments:
%      \cite[p.~4]{key}          \cite[see][p.~4]{key}
%
%  ---------------------------------------------------------------------------
%  KNOBS -- put any of these in main.tex AFTER \input{bib}
%  ---------------------------------------------------------------------------
%      \deckcitefootsize{\scriptsize}  larger bottom text   (default \tiny)
%      \deckcitefootrule{0pt}          drop the hairline above the block
%      \deckcitefootmax{6}             raise the crowding warning threshold
%      \deckcitefootoff                stop adding bottom references
%      \deckcitefooton                 turn them back on
%
%  \deckcitefootoff and \deckcitefooton also work INSIDE a frame, where they
%  apply to that frame only -- useful on a related-work slide with six cites.
%
%  ---------------------------------------------------------------------------
%  WHERE NOT TO PUT \cite
%  ---------------------------------------------------------------------------
%  Not in \frametitle, \section or \caption. Those are moving arguments: the
%  text is written out to the .toc/.aux and typeset again elsewhere, and a
%  footnote cannot survive that. Use \citequiet there.
% ============================================================================

\usepackage[
  backend     = biber,
  style       = numeric,   % [1] in the text, [1] in the reference list
  sorting     = none,      % list ordered by first citation, so [1] really is
                           % the first work you cite -- numbers read in order
  maxbibnames = 6,
  doi         = false,     % slides rarely need these; they eat space
  isbn        = false,
  url         = false,
  eprint      = false,
]{biblatex}

\usepackage{etoolbox}      % \forcsvlist -- biblatex loads it too, this line
                           % is here so the file does not rely on that

% ---------------------------------------------------------------------------
%  THE REFERENCES SLIDE
% ---------------------------------------------------------------------------
%  beamer replaces biblatex's entry label with its own "bibliography item"
%  template. Left empty it drops beamer's generic article icon -- but it also
%  silently deletes the [1] that numeric style depends on. \insertbiblabel
%  puts the number back.
\setbeamertemplate{bibliography item}{\insertbiblabel}
\renewcommand*{\bibfont}{\scriptsize}

% Keep entries from being hyphenated into a mess in a narrow column
\appto{\bibsetup}{\raggedright}

%  A frame that breaks across pages is titled "References I", "References II"
%  and so on. Numbering a list that turned out to fit on one slide just looks
%  like a mistake, so the numeral starts from the second page.
\setbeamertemplate{frametitle continuation}[from second]

% ---------------------------------------------------------------------------
%  A references slide that splits across pages automatically if it is long.
%  Use \bibliographyframe in your deck instead of writing this out by hand.
% ---------------------------------------------------------------------------
\newcommand{\bibliographyframe}[1][References]{%
  \begin{frame}[allowframebreaks]{#1}
    \printbibliography[heading=none]
  \end{frame}%
}

% ===========================================================================
%  THE PER-SLIDE REFERENCE BLOCK
% ===========================================================================
\makeatletter

% --- Knobs -----------------------------------------------------------------
\newcommand*{\deck@citefootsize}{\tiny}
\newcommand*{\deckcitefootsize}[1]{\renewcommand*{\deck@citefootsize}{#1}}

\newlength{\deck@citefootrule}
\setlength{\deck@citefootrule}{0.4pt}
\newcommand*{\deckcitefootrule}[1]{\setlength{\deck@citefootrule}{#1}}

\newcount\deck@citefootmax  \deck@citefootmax=4
\newcommand*{\deckcitefootmax}[1]{\deck@citefootmax=#1\relax}

\newif\ifdeck@citefoot  \deck@citefoottrue
\newcommand*{\deckcitefooton}{\deck@citefoottrue}
\newcommand*{\deckcitefootoff}{\deck@citefootfalse}

% --- The number, on its own ------------------------------------------------
%  Prints just [1] for a key, so the bottom line can be labelled with the
%  same number the reader just saw in the body text.
\DeclareCiteCommand{\deckcitenum}[\mkbibbrackets]
  {}
  {\printfield{labelnumber}}
  {\multicitedelim}
  {}

% --- The hairline above the block ------------------------------------------
%  This is LaTeX's footnote rule, so an ordinary \footnote gets the same line.
%  \deckcitefootrule{0pt} makes it disappear without disturbing the spacing.
\renewcommand{\footnoterule}{%
  \kern-3pt%
  \hbox{\color{BrandRule}%
        \vrule width 0.35\textwidth height \deck@citefootrule depth 0pt}%
  \kern2.6pt%
}

% --- One bottom line -------------------------------------------------------
%  An unmarked footnote: beamer already places footnotes at the bottom of the
%  frame, above the footer, and keeps them working under [allowframebreaks].
%  The mark is emptied and the footnote counter rolled back, because the line
%  carries its own [n] instead.
%
%  The footnote template is overridden LOCALLY. beamer's default reserves
%  1.8em for the mark; with no mark that is 1.8em of dead indent on every
%  line. Ordinary footnotes elsewhere in the deck are untouched.
%
%  The four \setbeamerfont lines are not optional. \fullcite goes through
%  beamer's bibliography drivers, which apply "bibliography entry ..." fonts --
%  absolute sizes that would override \deckcitefootsize and leave the [n]
%  label tiny next to a much larger reference. Unstarred \setbeamerfont
%  changes only the size and leaves series and shape alone, and all four are
%  local to the group, so the references slide at the end is unaffected.
\newcommand{\deck@citefootline}[1]{%
  \begingroup
    \setbeamerfont{bibliography entry author}  {size=\deck@citefootsize}%
    \setbeamerfont{bibliography entry title}   {size=\deck@citefootsize}%
    \setbeamerfont{bibliography entry location}{size=\deck@citefootsize}%
    \setbeamerfont{bibliography entry note}    {size=\deck@citefootsize}%
    \setbeamertemplate{footnote}{%
      \parindent0pt\noindent\raggedright
      \insertfootnotetext\par
    }%
    \renewcommand{\thefootnote}{}%
    \footnote{{\deck@citefootsize\deckcitenum{#1}\ \fullcite{#1}}}%
    \addtocounter{footnote}{-1}%
  \endgroup
}

% --- Crowding warning ------------------------------------------------------
%  Counts bottom references on the current slide and says so in the log once
%  the slide carries more than \deckcitefootmax of them. A crushed slide
%  should never be a surprise you only discover in the PDF.
\newcount\deck@citefootn
\def\deck@citefootwhere{}
\newcommand{\deck@citefootcount}{%
  \edef\deck@citefoothere{\the\c@page.\the\beamer@slideinframe}%
  \ifx\deck@citefoothere\deck@citefootwhere
    \global\advance\deck@citefootn\@ne
  \else
    \global\deck@citefootn\@ne
    \global\let\deck@citefootwhere\deck@citefoothere
  \fi
  \ifnum\deck@citefootn>\deck@citefootmax\relax
    \PackageWarning{deck-bib}{%
      \the\deck@citefootn\space bottom references on slide
      \the\c@page\space -- more than \the\deck@citefootmax.\MessageBreak
      They will crowd the slide. Use \string\citequiet\space for the
      extras,\MessageBreak
      or \string\deckcitefootoff\space at the top of that frame%
    }%
  \fi
}

% --- Dedup, per slide ------------------------------------------------------
%  Cite the same work twice on one slide and it should appear once at the
%  bottom; cite it again on a later slide and it should appear there too.
%
%  The marker is keyed on the physical page AND on which overlay of the frame
%  is being built. beamer typesets a frame body once per overlay, so a flag
%  reset by a begin-frame hook would fire once and then suppress the block on
%  overlays 2..n -- the reference would vanish the moment you pressed the
%  clicker. Keying on the slide being built needs no hook at all.
\newcommand{\deck@citefootone}[1]{%
  \ifcsname deck@seen@\the\c@page @\the\beamer@slideinframe @#1\endcsname
  \else
    \expandafter\gdef
      \csname deck@seen@\the\c@page @\the\beamer@slideinframe @#1\endcsname{}%
    \deck@citefootcount
    \deck@citefootline{#1}%
  \fi
}

%  \cite{a,b} is legal, so split the list. \forcsvlist trims stray spaces.
\newcommand{\deck@citefoot}[1]{%
  \ifdeck@citefoot
    \forcsvlist{\deck@citefootone}{#1}%
  \fi
}

% --- \cite does both -------------------------------------------------------
%  Wrapped, not reimplemented, so biblatex keeps full control of what [1]
%  actually looks like.
%
%  Deferred to \AtBeginDocument because hyperref -- which beamer loads for us
%  -- patches \cite at the end of the preamble. Grabbing it any earlier would
%  wrap a version that is about to be replaced.
%  biblatex reads ONE optional argument as the postnote and TWO as prenote +
%  postnote, so the wrapper has to hand back exactly as many as it was given.
%  Empty defaults are used rather than -NoValue- markers: an empty argument is
%  something \detokenize can test for certainty, and there is nothing to pass
%  through a second command that might mistake the marker for real text.
\newcommand{\deck@ifempty}[1]{%
  \if\relax\detokenize{#1}\relax
    \expandafter\@firstoftwo
  \else
    \expandafter\@secondoftwo
  \fi
}
\newcommand{\deck@passcite}[3]{%
  \deck@ifempty{#1}%
    {\deck@ifempty{#2}{\deck@origcite{#3}}{\deck@origcite[][#2]{#3}}}%
    {\deck@ifempty{#2}{\deck@origcite[#1]{#3}}{\deck@origcite[#1][#2]{#3}}}%
}

\NewDocumentCommand{\citequiet}{O{} O{} m}{\deck@passcite{#1}{#2}{#3}}

\AtBeginDocument{%
  \let\deck@origcite\cite
  \RenewDocumentCommand{\cite}{s O{} O{} m}{%
    \deck@passcite{#2}{#3}{#4}%
    \IfBooleanF{#1}{\deck@citefoot{#4}}%
  }%
}

\makeatother

%  ---------------------------------------------------------------------------
%      \deckcitefootsize{\scriptsize}  larger bottom text   (default \tiny)
%      \deckcitefootrule{0pt}          drop the hairline above the block
%      \deckcitefootmax{6}             raise the crowding warning threshold
%      \deckcitefootoff                stop adding bottom references
%      \deckcitefooton                 turn them back on
%
%  \deckcitefootoff and \deckcitefooton also work INSIDE a frame, where they
%  apply to that frame only -- useful on a related-work slide with six cites.
%
%  ---------------------------------------------------------------------------
%  WHERE NOT TO PUT \cite
%  ---------------------------------------------------------------------------
%  Not in \frametitle, \section or \caption. Those are moving arguments: the
%  text is written out to the .toc/.aux and typeset again elsewhere, and a
%  footnote cannot survive that. Use \citequiet there.
% ============================================================================

\usepackage[
  backend     = biber,
  style       = numeric,   % [1] in the text, [1] in the reference list
  sorting     = none,      % list ordered by first citation, so [1] really is
                           % the first work you cite -- numbers read in order
  maxbibnames = 6,
  doi         = false,     % slides rarely need these; they eat space
  isbn        = false,
  url         = false,
  eprint      = false,
]{biblatex}

\usepackage{etoolbox}      % \forcsvlist -- biblatex loads it too, this line
                           % is here so the file does not rely on that

% ---------------------------------------------------------------------------
%  THE REFERENCES SLIDE
% ---------------------------------------------------------------------------
%  beamer replaces biblatex's entry label with its own "bibliography item"
%  template. Left empty it drops beamer's generic article icon -- but it also
%  silently deletes the [1] that numeric style depends on. \insertbiblabel
%  puts the number back.
\setbeamertemplate{bibliography item}{\insertbiblabel}
\renewcommand*{\bibfont}{\scriptsize}

% Keep entries from being hyphenated into a mess in a narrow column
\appto{\bibsetup}{\raggedright}

%  A frame that breaks across pages is titled "References I", "References II"
%  and so on. Numbering a list that turned out to fit on one slide just looks
%  like a mistake, so the numeral starts from the second page.
\setbeamertemplate{frametitle continuation}[from second]

% ---------------------------------------------------------------------------
%  A references slide that splits across pages automatically if it is long.
%  Use \bibliographyframe in your deck instead of writing this out by hand.
% ---------------------------------------------------------------------------
\newcommand{\bibliographyframe}[1][References]{%
  \begin{frame}[allowframebreaks]{#1}
    \printbibliography[heading=none]
  \end{frame}%
}

% ===========================================================================
%  THE PER-SLIDE REFERENCE BLOCK
% ===========================================================================
\makeatletter

% --- Knobs -----------------------------------------------------------------
\newcommand*{\deck@citefootsize}{\tiny}
\newcommand*{\deckcitefootsize}[1]{\renewcommand*{\deck@citefootsize}{#1}}

\newlength{\deck@citefootrule}
\setlength{\deck@citefootrule}{0.4pt}
\newcommand*{\deckcitefootrule}[1]{\setlength{\deck@citefootrule}{#1}}

\newcount\deck@citefootmax  \deck@citefootmax=4
\newcommand*{\deckcitefootmax}[1]{\deck@citefootmax=#1\relax}

\newif\ifdeck@citefoot  \deck@citefoottrue
\newcommand*{\deckcitefooton}{\deck@citefoottrue}
\newcommand*{\deckcitefootoff}{\deck@citefootfalse}

% --- The number, on its own ------------------------------------------------
%  Prints just [1] for a key, so the bottom line can be labelled with the
%  same number the reader just saw in the body text.
\DeclareCiteCommand{\deckcitenum}[\mkbibbrackets]
  {}
  {\printfield{labelnumber}}
  {\multicitedelim}
  {}

% --- The hairline above the block ------------------------------------------
%  This is LaTeX's footnote rule, so an ordinary \footnote gets the same line.
%  \deckcitefootrule{0pt} makes it disappear without disturbing the spacing.
\renewcommand{\footnoterule}{%
  \kern-3pt%
  \hbox{\color{BrandRule}%
        \vrule width 0.35\textwidth height \deck@citefootrule depth 0pt}%
  \kern2.6pt%
}

% --- One bottom line -------------------------------------------------------
%  An unmarked footnote: beamer already places footnotes at the bottom of the
%  frame, above the footer, and keeps them working under [allowframebreaks].
%  The mark is emptied and the footnote counter rolled back, because the line
%  carries its own [n] instead.
%
%  The footnote template is overridden LOCALLY. beamer's default reserves
%  1.8em for the mark; with no mark that is 1.8em of dead indent on every
%  line. Ordinary footnotes elsewhere in the deck are untouched.
%
%  The four \setbeamerfont lines are not optional. \fullcite goes through
%  beamer's bibliography drivers, which apply "bibliography entry ..." fonts --
%  absolute sizes that would override \deckcitefootsize and leave the [n]
%  label tiny next to a much larger reference. Unstarred \setbeamerfont
%  changes only the size and leaves series and shape alone, and all four are
%  local to the group, so the references slide at the end is unaffected.
\newcommand{\deck@citefootline}[1]{%
  \begingroup
    \setbeamerfont{bibliography entry author}  {size=\deck@citefootsize}%
    \setbeamerfont{bibliography entry title}   {size=\deck@citefootsize}%
    \setbeamerfont{bibliography entry location}{size=\deck@citefootsize}%
    \setbeamerfont{bibliography entry note}    {size=\deck@citefootsize}%
    \setbeamertemplate{footnote}{%
      \parindent0pt\noindent\raggedright
      \insertfootnotetext\par
    }%
    \renewcommand{\thefootnote}{}%
    \footnote{{\deck@citefootsize\deckcitenum{#1}\ \fullcite{#1}}}%
    \addtocounter{footnote}{-1}%
  \endgroup
}

% --- Crowding warning ------------------------------------------------------
%  Counts bottom references on the current slide and says so in the log once
%  the slide carries more than \deckcitefootmax of them. A crushed slide
%  should never be a surprise you only discover in the PDF.
\newcount\deck@citefootn
\def\deck@citefootwhere{}
\newcommand{\deck@citefootcount}{%
  \edef\deck@citefoothere{\the\c@page.\the\beamer@slideinframe}%
  \ifx\deck@citefoothere\deck@citefootwhere
    \global\advance\deck@citefootn\@ne
  \else
    \global\deck@citefootn\@ne
    \global\let\deck@citefootwhere\deck@citefoothere
  \fi
  \ifnum\deck@citefootn>\deck@citefootmax\relax
    \PackageWarning{deck-bib}{%
      \the\deck@citefootn\space bottom references on slide
      \the\c@page\space -- more than \the\deck@citefootmax.\MessageBreak
      They will crowd the slide. Use \string\citequiet\space for the
      extras,\MessageBreak
      or \string\deckcitefootoff\space at the top of that frame%
    }%
  \fi
}

% --- Dedup, per slide ------------------------------------------------------
%  Cite the same work twice on one slide and it should appear once at the
%  bottom; cite it again on a later slide and it should appear there too.
%
%  The marker is keyed on the physical page AND on which overlay of the frame
%  is being built. beamer typesets a frame body once per overlay, so a flag
%  reset by a begin-frame hook would fire once and then suppress the block on
%  overlays 2..n -- the reference would vanish the moment you pressed the
%  clicker. Keying on the slide being built needs no hook at all.
\newcommand{\deck@citefootone}[1]{%
  \ifcsname deck@seen@\the\c@page @\the\beamer@slideinframe @#1\endcsname
  \else
    \expandafter\gdef
      \csname deck@seen@\the\c@page @\the\beamer@slideinframe @#1\endcsname{}%
    \deck@citefootcount
    \deck@citefootline{#1}%
  \fi
}

%  \cite{a,b} is legal, so split the list. \forcsvlist trims stray spaces.
\newcommand{\deck@citefoot}[1]{%
  \ifdeck@citefoot
    \forcsvlist{\deck@citefootone}{#1}%
  \fi
}

% --- \cite does both -------------------------------------------------------
%  Wrapped, not reimplemented, so biblatex keeps full control of what [1]
%  actually looks like.
%
%  Deferred to \AtBeginDocument because hyperref -- which beamer loads for us
%  -- patches \cite at the end of the preamble. Grabbing it any earlier would
%  wrap a version that is about to be replaced.
%  biblatex reads ONE optional argument as the postnote and TWO as prenote +
%  postnote, so the wrapper has to hand back exactly as many as it was given.
%  Empty defaults are used rather than -NoValue- markers: an empty argument is
%  something \detokenize can test for certainty, and there is nothing to pass
%  through a second command that might mistake the marker for real text.
\newcommand{\deck@ifempty}[1]{%
  \if\relax\detokenize{#1}\relax
    \expandafter\@firstoftwo
  \else
    \expandafter\@secondoftwo
  \fi
}
\newcommand{\deck@passcite}[3]{%
  \deck@ifempty{#1}%
    {\deck@ifempty{#2}{\deck@origcite{#3}}{\deck@origcite[][#2]{#3}}}%
    {\deck@ifempty{#2}{\deck@origcite[#1]{#3}}{\deck@origcite[#1][#2]{#3}}}%
}

\NewDocumentCommand{\citequiet}{O{} O{} m}{\deck@passcite{#1}{#2}{#3}}

\AtBeginDocument{%
  \let\deck@origcite\cite
  \RenewDocumentCommand{\cite}{s O{} O{} m}{%
    \deck@passcite{#2}{#3}{#4}%
    \IfBooleanF{#1}{\deck@citefoot{#4}}%
  }%
}

\makeatother

%      \addbibresource{refs.bib}
%      ...
%      Shown by \cite{knuth1984}.   ->  "Shown by [1]."  +  [1] at slide bottom
%      \citequiet{knuth1984}        ->  [1] only, nothing at the bottom
%      \cite*{knuth1984}            ->  same as \citequiet, shorter to type
%      \bibliographyframe           ->  the references slide at the end
%
%  \cite takes biblatex's usual optional arguments:
%      \cite[p.~4]{key}          \cite[see][p.~4]{key}
%
%  ---------------------------------------------------------------------------
%  KNOBS -- put any of these in main.tex AFTER % ============================================================================
%  preamble/bib.tex  --  citations and reference list
% ----------------------------------------------------------------------------
%  Numeric citations. \cite{key} does TWO things:
%
%      1. prints  [1]  where you wrote it
%      2. prints the FULL reference, small, at the bottom of THAT slide
%
%  and the same work appears again, with the same number, on the references
%  slide at the end. One command, nothing to keep in sync by hand.
%
%  In your deck:
%      % ============================================================================
%  preamble/bib.tex  --  citations and reference list
% ----------------------------------------------------------------------------
%  Numeric citations. \cite{key} does TWO things:
%
%      1. prints  [1]  where you wrote it
%      2. prints the FULL reference, small, at the bottom of THAT slide
%
%  and the same work appears again, with the same number, on the references
%  slide at the end. One command, nothing to keep in sync by hand.
%
%  In your deck:
%      \input{bib}
%      \addbibresource{refs.bib}
%      ...
%      Shown by \cite{knuth1984}.   ->  "Shown by [1]."  +  [1] at slide bottom
%      \citequiet{knuth1984}        ->  [1] only, nothing at the bottom
%      \cite*{knuth1984}            ->  same as \citequiet, shorter to type
%      \bibliographyframe           ->  the references slide at the end
%
%  \cite takes biblatex's usual optional arguments:
%      \cite[p.~4]{key}          \cite[see][p.~4]{key}
%
%  ---------------------------------------------------------------------------
%  KNOBS -- put any of these in main.tex AFTER \input{bib}
%  ---------------------------------------------------------------------------
%      \deckcitefootsize{\scriptsize}  larger bottom text   (default \tiny)
%      \deckcitefootrule{0pt}          drop the hairline above the block
%      \deckcitefootmax{6}             raise the crowding warning threshold
%      \deckcitefootoff                stop adding bottom references
%      \deckcitefooton                 turn them back on
%
%  \deckcitefootoff and \deckcitefooton also work INSIDE a frame, where they
%  apply to that frame only -- useful on a related-work slide with six cites.
%
%  ---------------------------------------------------------------------------
%  WHERE NOT TO PUT \cite
%  ---------------------------------------------------------------------------
%  Not in \frametitle, \section or \caption. Those are moving arguments: the
%  text is written out to the .toc/.aux and typeset again elsewhere, and a
%  footnote cannot survive that. Use \citequiet there.
% ============================================================================

\usepackage[
  backend     = biber,
  style       = numeric,   % [1] in the text, [1] in the reference list
  sorting     = none,      % list ordered by first citation, so [1] really is
                           % the first work you cite -- numbers read in order
  maxbibnames = 6,
  doi         = false,     % slides rarely need these; they eat space
  isbn        = false,
  url         = false,
  eprint      = false,
]{biblatex}

\usepackage{etoolbox}      % \forcsvlist -- biblatex loads it too, this line
                           % is here so the file does not rely on that

% ---------------------------------------------------------------------------
%  THE REFERENCES SLIDE
% ---------------------------------------------------------------------------
%  beamer replaces biblatex's entry label with its own "bibliography item"
%  template. Left empty it drops beamer's generic article icon -- but it also
%  silently deletes the [1] that numeric style depends on. \insertbiblabel
%  puts the number back.
\setbeamertemplate{bibliography item}{\insertbiblabel}
\renewcommand*{\bibfont}{\scriptsize}

% Keep entries from being hyphenated into a mess in a narrow column
\appto{\bibsetup}{\raggedright}

%  A frame that breaks across pages is titled "References I", "References II"
%  and so on. Numbering a list that turned out to fit on one slide just looks
%  like a mistake, so the numeral starts from the second page.
\setbeamertemplate{frametitle continuation}[from second]

% ---------------------------------------------------------------------------
%  A references slide that splits across pages automatically if it is long.
%  Use \bibliographyframe in your deck instead of writing this out by hand.
% ---------------------------------------------------------------------------
\newcommand{\bibliographyframe}[1][References]{%
  \begin{frame}[allowframebreaks]{#1}
    \printbibliography[heading=none]
  \end{frame}%
}

% ===========================================================================
%  THE PER-SLIDE REFERENCE BLOCK
% ===========================================================================
\makeatletter

% --- Knobs -----------------------------------------------------------------
\newcommand*{\deck@citefootsize}{\tiny}
\newcommand*{\deckcitefootsize}[1]{\renewcommand*{\deck@citefootsize}{#1}}

\newlength{\deck@citefootrule}
\setlength{\deck@citefootrule}{0.4pt}
\newcommand*{\deckcitefootrule}[1]{\setlength{\deck@citefootrule}{#1}}

\newcount\deck@citefootmax  \deck@citefootmax=4
\newcommand*{\deckcitefootmax}[1]{\deck@citefootmax=#1\relax}

\newif\ifdeck@citefoot  \deck@citefoottrue
\newcommand*{\deckcitefooton}{\deck@citefoottrue}
\newcommand*{\deckcitefootoff}{\deck@citefootfalse}

% --- The number, on its own ------------------------------------------------
%  Prints just [1] for a key, so the bottom line can be labelled with the
%  same number the reader just saw in the body text.
\DeclareCiteCommand{\deckcitenum}[\mkbibbrackets]
  {}
  {\printfield{labelnumber}}
  {\multicitedelim}
  {}

% --- The hairline above the block ------------------------------------------
%  This is LaTeX's footnote rule, so an ordinary \footnote gets the same line.
%  \deckcitefootrule{0pt} makes it disappear without disturbing the spacing.
\renewcommand{\footnoterule}{%
  \kern-3pt%
  \hbox{\color{BrandRule}%
        \vrule width 0.35\textwidth height \deck@citefootrule depth 0pt}%
  \kern2.6pt%
}

% --- One bottom line -------------------------------------------------------
%  An unmarked footnote: beamer already places footnotes at the bottom of the
%  frame, above the footer, and keeps them working under [allowframebreaks].
%  The mark is emptied and the footnote counter rolled back, because the line
%  carries its own [n] instead.
%
%  The footnote template is overridden LOCALLY. beamer's default reserves
%  1.8em for the mark; with no mark that is 1.8em of dead indent on every
%  line. Ordinary footnotes elsewhere in the deck are untouched.
%
%  The four \setbeamerfont lines are not optional. \fullcite goes through
%  beamer's bibliography drivers, which apply "bibliography entry ..." fonts --
%  absolute sizes that would override \deckcitefootsize and leave the [n]
%  label tiny next to a much larger reference. Unstarred \setbeamerfont
%  changes only the size and leaves series and shape alone, and all four are
%  local to the group, so the references slide at the end is unaffected.
\newcommand{\deck@citefootline}[1]{%
  \begingroup
    \setbeamerfont{bibliography entry author}  {size=\deck@citefootsize}%
    \setbeamerfont{bibliography entry title}   {size=\deck@citefootsize}%
    \setbeamerfont{bibliography entry location}{size=\deck@citefootsize}%
    \setbeamerfont{bibliography entry note}    {size=\deck@citefootsize}%
    \setbeamertemplate{footnote}{%
      \parindent0pt\noindent\raggedright
      \insertfootnotetext\par
    }%
    \renewcommand{\thefootnote}{}%
    \footnote{{\deck@citefootsize\deckcitenum{#1}\ \fullcite{#1}}}%
    \addtocounter{footnote}{-1}%
  \endgroup
}

% --- Crowding warning ------------------------------------------------------
%  Counts bottom references on the current slide and says so in the log once
%  the slide carries more than \deckcitefootmax of them. A crushed slide
%  should never be a surprise you only discover in the PDF.
\newcount\deck@citefootn
\def\deck@citefootwhere{}
\newcommand{\deck@citefootcount}{%
  \edef\deck@citefoothere{\the\c@page.\the\beamer@slideinframe}%
  \ifx\deck@citefoothere\deck@citefootwhere
    \global\advance\deck@citefootn\@ne
  \else
    \global\deck@citefootn\@ne
    \global\let\deck@citefootwhere\deck@citefoothere
  \fi
  \ifnum\deck@citefootn>\deck@citefootmax\relax
    \PackageWarning{deck-bib}{%
      \the\deck@citefootn\space bottom references on slide
      \the\c@page\space -- more than \the\deck@citefootmax.\MessageBreak
      They will crowd the slide. Use \string\citequiet\space for the
      extras,\MessageBreak
      or \string\deckcitefootoff\space at the top of that frame%
    }%
  \fi
}

% --- Dedup, per slide ------------------------------------------------------
%  Cite the same work twice on one slide and it should appear once at the
%  bottom; cite it again on a later slide and it should appear there too.
%
%  The marker is keyed on the physical page AND on which overlay of the frame
%  is being built. beamer typesets a frame body once per overlay, so a flag
%  reset by a begin-frame hook would fire once and then suppress the block on
%  overlays 2..n -- the reference would vanish the moment you pressed the
%  clicker. Keying on the slide being built needs no hook at all.
\newcommand{\deck@citefootone}[1]{%
  \ifcsname deck@seen@\the\c@page @\the\beamer@slideinframe @#1\endcsname
  \else
    \expandafter\gdef
      \csname deck@seen@\the\c@page @\the\beamer@slideinframe @#1\endcsname{}%
    \deck@citefootcount
    \deck@citefootline{#1}%
  \fi
}

%  \cite{a,b} is legal, so split the list. \forcsvlist trims stray spaces.
\newcommand{\deck@citefoot}[1]{%
  \ifdeck@citefoot
    \forcsvlist{\deck@citefootone}{#1}%
  \fi
}

% --- \cite does both -------------------------------------------------------
%  Wrapped, not reimplemented, so biblatex keeps full control of what [1]
%  actually looks like.
%
%  Deferred to \AtBeginDocument because hyperref -- which beamer loads for us
%  -- patches \cite at the end of the preamble. Grabbing it any earlier would
%  wrap a version that is about to be replaced.
%  biblatex reads ONE optional argument as the postnote and TWO as prenote +
%  postnote, so the wrapper has to hand back exactly as many as it was given.
%  Empty defaults are used rather than -NoValue- markers: an empty argument is
%  something \detokenize can test for certainty, and there is nothing to pass
%  through a second command that might mistake the marker for real text.
\newcommand{\deck@ifempty}[1]{%
  \if\relax\detokenize{#1}\relax
    \expandafter\@firstoftwo
  \else
    \expandafter\@secondoftwo
  \fi
}
\newcommand{\deck@passcite}[3]{%
  \deck@ifempty{#1}%
    {\deck@ifempty{#2}{\deck@origcite{#3}}{\deck@origcite[][#2]{#3}}}%
    {\deck@ifempty{#2}{\deck@origcite[#1]{#3}}{\deck@origcite[#1][#2]{#3}}}%
}

\NewDocumentCommand{\citequiet}{O{} O{} m}{\deck@passcite{#1}{#2}{#3}}

\AtBeginDocument{%
  \let\deck@origcite\cite
  \RenewDocumentCommand{\cite}{s O{} O{} m}{%
    \deck@passcite{#2}{#3}{#4}%
    \IfBooleanF{#1}{\deck@citefoot{#4}}%
  }%
}

\makeatother

%      \addbibresource{refs.bib}
%      ...
%      Shown by \cite{knuth1984}.   ->  "Shown by [1]."  +  [1] at slide bottom
%      \citequiet{knuth1984}        ->  [1] only, nothing at the bottom
%      \cite*{knuth1984}            ->  same as \citequiet, shorter to type
%      \bibliographyframe           ->  the references slide at the end
%
%  \cite takes biblatex's usual optional arguments:
%      \cite[p.~4]{key}          \cite[see][p.~4]{key}
%
%  ---------------------------------------------------------------------------
%  KNOBS -- put any of these in main.tex AFTER % ============================================================================
%  preamble/bib.tex  --  citations and reference list
% ----------------------------------------------------------------------------
%  Numeric citations. \cite{key} does TWO things:
%
%      1. prints  [1]  where you wrote it
%      2. prints the FULL reference, small, at the bottom of THAT slide
%
%  and the same work appears again, with the same number, on the references
%  slide at the end. One command, nothing to keep in sync by hand.
%
%  In your deck:
%      \input{bib}
%      \addbibresource{refs.bib}
%      ...
%      Shown by \cite{knuth1984}.   ->  "Shown by [1]."  +  [1] at slide bottom
%      \citequiet{knuth1984}        ->  [1] only, nothing at the bottom
%      \cite*{knuth1984}            ->  same as \citequiet, shorter to type
%      \bibliographyframe           ->  the references slide at the end
%
%  \cite takes biblatex's usual optional arguments:
%      \cite[p.~4]{key}          \cite[see][p.~4]{key}
%
%  ---------------------------------------------------------------------------
%  KNOBS -- put any of these in main.tex AFTER \input{bib}
%  ---------------------------------------------------------------------------
%      \deckcitefootsize{\scriptsize}  larger bottom text   (default \tiny)
%      \deckcitefootrule{0pt}          drop the hairline above the block
%      \deckcitefootmax{6}             raise the crowding warning threshold
%      \deckcitefootoff                stop adding bottom references
%      \deckcitefooton                 turn them back on
%
%  \deckcitefootoff and \deckcitefooton also work INSIDE a frame, where they
%  apply to that frame only -- useful on a related-work slide with six cites.
%
%  ---------------------------------------------------------------------------
%  WHERE NOT TO PUT \cite
%  ---------------------------------------------------------------------------
%  Not in \frametitle, \section or \caption. Those are moving arguments: the
%  text is written out to the .toc/.aux and typeset again elsewhere, and a
%  footnote cannot survive that. Use \citequiet there.
% ============================================================================

\usepackage[
  backend     = biber,
  style       = numeric,   % [1] in the text, [1] in the reference list
  sorting     = none,      % list ordered by first citation, so [1] really is
                           % the first work you cite -- numbers read in order
  maxbibnames = 6,
  doi         = false,     % slides rarely need these; they eat space
  isbn        = false,
  url         = false,
  eprint      = false,
]{biblatex}

\usepackage{etoolbox}      % \forcsvlist -- biblatex loads it too, this line
                           % is here so the file does not rely on that

% ---------------------------------------------------------------------------
%  THE REFERENCES SLIDE
% ---------------------------------------------------------------------------
%  beamer replaces biblatex's entry label with its own "bibliography item"
%  template. Left empty it drops beamer's generic article icon -- but it also
%  silently deletes the [1] that numeric style depends on. \insertbiblabel
%  puts the number back.
\setbeamertemplate{bibliography item}{\insertbiblabel}
\renewcommand*{\bibfont}{\scriptsize}

% Keep entries from being hyphenated into a mess in a narrow column
\appto{\bibsetup}{\raggedright}

%  A frame that breaks across pages is titled "References I", "References II"
%  and so on. Numbering a list that turned out to fit on one slide just looks
%  like a mistake, so the numeral starts from the second page.
\setbeamertemplate{frametitle continuation}[from second]

% ---------------------------------------------------------------------------
%  A references slide that splits across pages automatically if it is long.
%  Use \bibliographyframe in your deck instead of writing this out by hand.
% ---------------------------------------------------------------------------
\newcommand{\bibliographyframe}[1][References]{%
  \begin{frame}[allowframebreaks]{#1}
    \printbibliography[heading=none]
  \end{frame}%
}

% ===========================================================================
%  THE PER-SLIDE REFERENCE BLOCK
% ===========================================================================
\makeatletter

% --- Knobs -----------------------------------------------------------------
\newcommand*{\deck@citefootsize}{\tiny}
\newcommand*{\deckcitefootsize}[1]{\renewcommand*{\deck@citefootsize}{#1}}

\newlength{\deck@citefootrule}
\setlength{\deck@citefootrule}{0.4pt}
\newcommand*{\deckcitefootrule}[1]{\setlength{\deck@citefootrule}{#1}}

\newcount\deck@citefootmax  \deck@citefootmax=4
\newcommand*{\deckcitefootmax}[1]{\deck@citefootmax=#1\relax}

\newif\ifdeck@citefoot  \deck@citefoottrue
\newcommand*{\deckcitefooton}{\deck@citefoottrue}
\newcommand*{\deckcitefootoff}{\deck@citefootfalse}

% --- The number, on its own ------------------------------------------------
%  Prints just [1] for a key, so the bottom line can be labelled with the
%  same number the reader just saw in the body text.
\DeclareCiteCommand{\deckcitenum}[\mkbibbrackets]
  {}
  {\printfield{labelnumber}}
  {\multicitedelim}
  {}

% --- The hairline above the block ------------------------------------------
%  This is LaTeX's footnote rule, so an ordinary \footnote gets the same line.
%  \deckcitefootrule{0pt} makes it disappear without disturbing the spacing.
\renewcommand{\footnoterule}{%
  \kern-3pt%
  \hbox{\color{BrandRule}%
        \vrule width 0.35\textwidth height \deck@citefootrule depth 0pt}%
  \kern2.6pt%
}

% --- One bottom line -------------------------------------------------------
%  An unmarked footnote: beamer already places footnotes at the bottom of the
%  frame, above the footer, and keeps them working under [allowframebreaks].
%  The mark is emptied and the footnote counter rolled back, because the line
%  carries its own [n] instead.
%
%  The footnote template is overridden LOCALLY. beamer's default reserves
%  1.8em for the mark; with no mark that is 1.8em of dead indent on every
%  line. Ordinary footnotes elsewhere in the deck are untouched.
%
%  The four \setbeamerfont lines are not optional. \fullcite goes through
%  beamer's bibliography drivers, which apply "bibliography entry ..." fonts --
%  absolute sizes that would override \deckcitefootsize and leave the [n]
%  label tiny next to a much larger reference. Unstarred \setbeamerfont
%  changes only the size and leaves series and shape alone, and all four are
%  local to the group, so the references slide at the end is unaffected.
\newcommand{\deck@citefootline}[1]{%
  \begingroup
    \setbeamerfont{bibliography entry author}  {size=\deck@citefootsize}%
    \setbeamerfont{bibliography entry title}   {size=\deck@citefootsize}%
    \setbeamerfont{bibliography entry location}{size=\deck@citefootsize}%
    \setbeamerfont{bibliography entry note}    {size=\deck@citefootsize}%
    \setbeamertemplate{footnote}{%
      \parindent0pt\noindent\raggedright
      \insertfootnotetext\par
    }%
    \renewcommand{\thefootnote}{}%
    \footnote{{\deck@citefootsize\deckcitenum{#1}\ \fullcite{#1}}}%
    \addtocounter{footnote}{-1}%
  \endgroup
}

% --- Crowding warning ------------------------------------------------------
%  Counts bottom references on the current slide and says so in the log once
%  the slide carries more than \deckcitefootmax of them. A crushed slide
%  should never be a surprise you only discover in the PDF.
\newcount\deck@citefootn
\def\deck@citefootwhere{}
\newcommand{\deck@citefootcount}{%
  \edef\deck@citefoothere{\the\c@page.\the\beamer@slideinframe}%
  \ifx\deck@citefoothere\deck@citefootwhere
    \global\advance\deck@citefootn\@ne
  \else
    \global\deck@citefootn\@ne
    \global\let\deck@citefootwhere\deck@citefoothere
  \fi
  \ifnum\deck@citefootn>\deck@citefootmax\relax
    \PackageWarning{deck-bib}{%
      \the\deck@citefootn\space bottom references on slide
      \the\c@page\space -- more than \the\deck@citefootmax.\MessageBreak
      They will crowd the slide. Use \string\citequiet\space for the
      extras,\MessageBreak
      or \string\deckcitefootoff\space at the top of that frame%
    }%
  \fi
}

% --- Dedup, per slide ------------------------------------------------------
%  Cite the same work twice on one slide and it should appear once at the
%  bottom; cite it again on a later slide and it should appear there too.
%
%  The marker is keyed on the physical page AND on which overlay of the frame
%  is being built. beamer typesets a frame body once per overlay, so a flag
%  reset by a begin-frame hook would fire once and then suppress the block on
%  overlays 2..n -- the reference would vanish the moment you pressed the
%  clicker. Keying on the slide being built needs no hook at all.
\newcommand{\deck@citefootone}[1]{%
  \ifcsname deck@seen@\the\c@page @\the\beamer@slideinframe @#1\endcsname
  \else
    \expandafter\gdef
      \csname deck@seen@\the\c@page @\the\beamer@slideinframe @#1\endcsname{}%
    \deck@citefootcount
    \deck@citefootline{#1}%
  \fi
}

%  \cite{a,b} is legal, so split the list. \forcsvlist trims stray spaces.
\newcommand{\deck@citefoot}[1]{%
  \ifdeck@citefoot
    \forcsvlist{\deck@citefootone}{#1}%
  \fi
}

% --- \cite does both -------------------------------------------------------
%  Wrapped, not reimplemented, so biblatex keeps full control of what [1]
%  actually looks like.
%
%  Deferred to \AtBeginDocument because hyperref -- which beamer loads for us
%  -- patches \cite at the end of the preamble. Grabbing it any earlier would
%  wrap a version that is about to be replaced.
%  biblatex reads ONE optional argument as the postnote and TWO as prenote +
%  postnote, so the wrapper has to hand back exactly as many as it was given.
%  Empty defaults are used rather than -NoValue- markers: an empty argument is
%  something \detokenize can test for certainty, and there is nothing to pass
%  through a second command that might mistake the marker for real text.
\newcommand{\deck@ifempty}[1]{%
  \if\relax\detokenize{#1}\relax
    \expandafter\@firstoftwo
  \else
    \expandafter\@secondoftwo
  \fi
}
\newcommand{\deck@passcite}[3]{%
  \deck@ifempty{#1}%
    {\deck@ifempty{#2}{\deck@origcite{#3}}{\deck@origcite[][#2]{#3}}}%
    {\deck@ifempty{#2}{\deck@origcite[#1]{#3}}{\deck@origcite[#1][#2]{#3}}}%
}

\NewDocumentCommand{\citequiet}{O{} O{} m}{\deck@passcite{#1}{#2}{#3}}

\AtBeginDocument{%
  \let\deck@origcite\cite
  \RenewDocumentCommand{\cite}{s O{} O{} m}{%
    \deck@passcite{#2}{#3}{#4}%
    \IfBooleanF{#1}{\deck@citefoot{#4}}%
  }%
}

\makeatother

%  ---------------------------------------------------------------------------
%      \deckcitefootsize{\scriptsize}  larger bottom text   (default \tiny)
%      \deckcitefootrule{0pt}          drop the hairline above the block
%      \deckcitefootmax{6}             raise the crowding warning threshold
%      \deckcitefootoff                stop adding bottom references
%      \deckcitefooton                 turn them back on
%
%  \deckcitefootoff and \deckcitefooton also work INSIDE a frame, where they
%  apply to that frame only -- useful on a related-work slide with six cites.
%
%  ---------------------------------------------------------------------------
%  WHERE NOT TO PUT \cite
%  ---------------------------------------------------------------------------
%  Not in \frametitle, \section or \caption. Those are moving arguments: the
%  text is written out to the .toc/.aux and typeset again elsewhere, and a
%  footnote cannot survive that. Use \citequiet there.
% ============================================================================

\usepackage[
  backend     = biber,
  style       = numeric,   % [1] in the text, [1] in the reference list
  sorting     = none,      % list ordered by first citation, so [1] really is
                           % the first work you cite -- numbers read in order
  maxbibnames = 6,
  doi         = false,     % slides rarely need these; they eat space
  isbn        = false,
  url         = false,
  eprint      = false,
]{biblatex}

\usepackage{etoolbox}      % \forcsvlist -- biblatex loads it too, this line
                           % is here so the file does not rely on that

% ---------------------------------------------------------------------------
%  THE REFERENCES SLIDE
% ---------------------------------------------------------------------------
%  beamer replaces biblatex's entry label with its own "bibliography item"
%  template. Left empty it drops beamer's generic article icon -- but it also
%  silently deletes the [1] that numeric style depends on. \insertbiblabel
%  puts the number back.
\setbeamertemplate{bibliography item}{\insertbiblabel}
\renewcommand*{\bibfont}{\scriptsize}

% Keep entries from being hyphenated into a mess in a narrow column
\appto{\bibsetup}{\raggedright}

%  A frame that breaks across pages is titled "References I", "References II"
%  and so on. Numbering a list that turned out to fit on one slide just looks
%  like a mistake, so the numeral starts from the second page.
\setbeamertemplate{frametitle continuation}[from second]

% ---------------------------------------------------------------------------
%  A references slide that splits across pages automatically if it is long.
%  Use \bibliographyframe in your deck instead of writing this out by hand.
% ---------------------------------------------------------------------------
\newcommand{\bibliographyframe}[1][References]{%
  \begin{frame}[allowframebreaks]{#1}
    \printbibliography[heading=none]
  \end{frame}%
}

% ===========================================================================
%  THE PER-SLIDE REFERENCE BLOCK
% ===========================================================================
\makeatletter

% --- Knobs -----------------------------------------------------------------
\newcommand*{\deck@citefootsize}{\tiny}
\newcommand*{\deckcitefootsize}[1]{\renewcommand*{\deck@citefootsize}{#1}}

\newlength{\deck@citefootrule}
\setlength{\deck@citefootrule}{0.4pt}
\newcommand*{\deckcitefootrule}[1]{\setlength{\deck@citefootrule}{#1}}

\newcount\deck@citefootmax  \deck@citefootmax=4
\newcommand*{\deckcitefootmax}[1]{\deck@citefootmax=#1\relax}

\newif\ifdeck@citefoot  \deck@citefoottrue
\newcommand*{\deckcitefooton}{\deck@citefoottrue}
\newcommand*{\deckcitefootoff}{\deck@citefootfalse}

% --- The number, on its own ------------------------------------------------
%  Prints just [1] for a key, so the bottom line can be labelled with the
%  same number the reader just saw in the body text.
\DeclareCiteCommand{\deckcitenum}[\mkbibbrackets]
  {}
  {\printfield{labelnumber}}
  {\multicitedelim}
  {}

% --- The hairline above the block ------------------------------------------
%  This is LaTeX's footnote rule, so an ordinary \footnote gets the same line.
%  \deckcitefootrule{0pt} makes it disappear without disturbing the spacing.
\renewcommand{\footnoterule}{%
  \kern-3pt%
  \hbox{\color{BrandRule}%
        \vrule width 0.35\textwidth height \deck@citefootrule depth 0pt}%
  \kern2.6pt%
}

% --- One bottom line -------------------------------------------------------
%  An unmarked footnote: beamer already places footnotes at the bottom of the
%  frame, above the footer, and keeps them working under [allowframebreaks].
%  The mark is emptied and the footnote counter rolled back, because the line
%  carries its own [n] instead.
%
%  The footnote template is overridden LOCALLY. beamer's default reserves
%  1.8em for the mark; with no mark that is 1.8em of dead indent on every
%  line. Ordinary footnotes elsewhere in the deck are untouched.
%
%  The four \setbeamerfont lines are not optional. \fullcite goes through
%  beamer's bibliography drivers, which apply "bibliography entry ..." fonts --
%  absolute sizes that would override \deckcitefootsize and leave the [n]
%  label tiny next to a much larger reference. Unstarred \setbeamerfont
%  changes only the size and leaves series and shape alone, and all four are
%  local to the group, so the references slide at the end is unaffected.
\newcommand{\deck@citefootline}[1]{%
  \begingroup
    \setbeamerfont{bibliography entry author}  {size=\deck@citefootsize}%
    \setbeamerfont{bibliography entry title}   {size=\deck@citefootsize}%
    \setbeamerfont{bibliography entry location}{size=\deck@citefootsize}%
    \setbeamerfont{bibliography entry note}    {size=\deck@citefootsize}%
    \setbeamertemplate{footnote}{%
      \parindent0pt\noindent\raggedright
      \insertfootnotetext\par
    }%
    \renewcommand{\thefootnote}{}%
    \footnote{{\deck@citefootsize\deckcitenum{#1}\ \fullcite{#1}}}%
    \addtocounter{footnote}{-1}%
  \endgroup
}

% --- Crowding warning ------------------------------------------------------
%  Counts bottom references on the current slide and says so in the log once
%  the slide carries more than \deckcitefootmax of them. A crushed slide
%  should never be a surprise you only discover in the PDF.
\newcount\deck@citefootn
\def\deck@citefootwhere{}
\newcommand{\deck@citefootcount}{%
  \edef\deck@citefoothere{\the\c@page.\the\beamer@slideinframe}%
  \ifx\deck@citefoothere\deck@citefootwhere
    \global\advance\deck@citefootn\@ne
  \else
    \global\deck@citefootn\@ne
    \global\let\deck@citefootwhere\deck@citefoothere
  \fi
  \ifnum\deck@citefootn>\deck@citefootmax\relax
    \PackageWarning{deck-bib}{%
      \the\deck@citefootn\space bottom references on slide
      \the\c@page\space -- more than \the\deck@citefootmax.\MessageBreak
      They will crowd the slide. Use \string\citequiet\space for the
      extras,\MessageBreak
      or \string\deckcitefootoff\space at the top of that frame%
    }%
  \fi
}

% --- Dedup, per slide ------------------------------------------------------
%  Cite the same work twice on one slide and it should appear once at the
%  bottom; cite it again on a later slide and it should appear there too.
%
%  The marker is keyed on the physical page AND on which overlay of the frame
%  is being built. beamer typesets a frame body once per overlay, so a flag
%  reset by a begin-frame hook would fire once and then suppress the block on
%  overlays 2..n -- the reference would vanish the moment you pressed the
%  clicker. Keying on the slide being built needs no hook at all.
\newcommand{\deck@citefootone}[1]{%
  \ifcsname deck@seen@\the\c@page @\the\beamer@slideinframe @#1\endcsname
  \else
    \expandafter\gdef
      \csname deck@seen@\the\c@page @\the\beamer@slideinframe @#1\endcsname{}%
    \deck@citefootcount
    \deck@citefootline{#1}%
  \fi
}

%  \cite{a,b} is legal, so split the list. \forcsvlist trims stray spaces.
\newcommand{\deck@citefoot}[1]{%
  \ifdeck@citefoot
    \forcsvlist{\deck@citefootone}{#1}%
  \fi
}

% --- \cite does both -------------------------------------------------------
%  Wrapped, not reimplemented, so biblatex keeps full control of what [1]
%  actually looks like.
%
%  Deferred to \AtBeginDocument because hyperref -- which beamer loads for us
%  -- patches \cite at the end of the preamble. Grabbing it any earlier would
%  wrap a version that is about to be replaced.
%  biblatex reads ONE optional argument as the postnote and TWO as prenote +
%  postnote, so the wrapper has to hand back exactly as many as it was given.
%  Empty defaults are used rather than -NoValue- markers: an empty argument is
%  something \detokenize can test for certainty, and there is nothing to pass
%  through a second command that might mistake the marker for real text.
\newcommand{\deck@ifempty}[1]{%
  \if\relax\detokenize{#1}\relax
    \expandafter\@firstoftwo
  \else
    \expandafter\@secondoftwo
  \fi
}
\newcommand{\deck@passcite}[3]{%
  \deck@ifempty{#1}%
    {\deck@ifempty{#2}{\deck@origcite{#3}}{\deck@origcite[][#2]{#3}}}%
    {\deck@ifempty{#2}{\deck@origcite[#1]{#3}}{\deck@origcite[#1][#2]{#3}}}%
}

\NewDocumentCommand{\citequiet}{O{} O{} m}{\deck@passcite{#1}{#2}{#3}}

\AtBeginDocument{%
  \let\deck@origcite\cite
  \RenewDocumentCommand{\cite}{s O{} O{} m}{%
    \deck@passcite{#2}{#3}{#4}%
    \IfBooleanF{#1}{\deck@citefoot{#4}}%
  }%
}

\makeatother

%  ---------------------------------------------------------------------------
%      \deckcitefootsize{\scriptsize}  larger bottom text   (default \tiny)
%      \deckcitefootrule{0pt}          drop the hairline above the block
%      \deckcitefootmax{6}             raise the crowding warning threshold
%      \deckcitefootoff                stop adding bottom references
%      \deckcitefooton                 turn them back on
%
%  \deckcitefootoff and \deckcitefooton also work INSIDE a frame, where they
%  apply to that frame only -- useful on a related-work slide with six cites.
%
%  ---------------------------------------------------------------------------
%  WHERE NOT TO PUT \cite
%  ---------------------------------------------------------------------------
%  Not in \frametitle, \section or \caption. Those are moving arguments: the
%  text is written out to the .toc/.aux and typeset again elsewhere, and a
%  footnote cannot survive that. Use \citequiet there.
% ============================================================================

\usepackage[
  backend     = biber,
  style       = numeric,   % [1] in the text, [1] in the reference list
  sorting     = none,      % list ordered by first citation, so [1] really is
                           % the first work you cite -- numbers read in order
  maxbibnames = 6,
  doi         = false,     % slides rarely need these; they eat space
  isbn        = false,
  url         = false,
  eprint      = false,
]{biblatex}

\usepackage{etoolbox}      % \forcsvlist -- biblatex loads it too, this line
                           % is here so the file does not rely on that

% ---------------------------------------------------------------------------
%  THE REFERENCES SLIDE
% ---------------------------------------------------------------------------
%  beamer replaces biblatex's entry label with its own "bibliography item"
%  template. Left empty it drops beamer's generic article icon -- but it also
%  silently deletes the [1] that numeric style depends on. \insertbiblabel
%  puts the number back.
\setbeamertemplate{bibliography item}{\insertbiblabel}
\renewcommand*{\bibfont}{\scriptsize}

% Keep entries from being hyphenated into a mess in a narrow column
\appto{\bibsetup}{\raggedright}

%  A frame that breaks across pages is titled "References I", "References II"
%  and so on. Numbering a list that turned out to fit on one slide just looks
%  like a mistake, so the numeral starts from the second page.
\setbeamertemplate{frametitle continuation}[from second]

% ---------------------------------------------------------------------------
%  A references slide that splits across pages automatically if it is long.
%  Use \bibliographyframe in your deck instead of writing this out by hand.
% ---------------------------------------------------------------------------
\newcommand{\bibliographyframe}[1][References]{%
  \begin{frame}[allowframebreaks]{#1}
    \printbibliography[heading=none]
  \end{frame}%
}

% ===========================================================================
%  THE PER-SLIDE REFERENCE BLOCK
% ===========================================================================
\makeatletter

% --- Knobs -----------------------------------------------------------------
\newcommand*{\deck@citefootsize}{\tiny}
\newcommand*{\deckcitefootsize}[1]{\renewcommand*{\deck@citefootsize}{#1}}

\newlength{\deck@citefootrule}
\setlength{\deck@citefootrule}{0.4pt}
\newcommand*{\deckcitefootrule}[1]{\setlength{\deck@citefootrule}{#1}}

\newcount\deck@citefootmax  \deck@citefootmax=4
\newcommand*{\deckcitefootmax}[1]{\deck@citefootmax=#1\relax}

\newif\ifdeck@citefoot  \deck@citefoottrue
\newcommand*{\deckcitefooton}{\deck@citefoottrue}
\newcommand*{\deckcitefootoff}{\deck@citefootfalse}

% --- The number, on its own ------------------------------------------------
%  Prints just [1] for a key, so the bottom line can be labelled with the
%  same number the reader just saw in the body text.
\DeclareCiteCommand{\deckcitenum}[\mkbibbrackets]
  {}
  {\printfield{labelnumber}}
  {\multicitedelim}
  {}

% --- The hairline above the block ------------------------------------------
%  This is LaTeX's footnote rule, so an ordinary \footnote gets the same line.
%  \deckcitefootrule{0pt} makes it disappear without disturbing the spacing.
\renewcommand{\footnoterule}{%
  \kern-3pt%
  \hbox{\color{BrandRule}%
        \vrule width 0.35\textwidth height \deck@citefootrule depth 0pt}%
  \kern2.6pt%
}

% --- One bottom line -------------------------------------------------------
%  An unmarked footnote: beamer already places footnotes at the bottom of the
%  frame, above the footer, and keeps them working under [allowframebreaks].
%  The mark is emptied and the footnote counter rolled back, because the line
%  carries its own [n] instead.
%
%  The footnote template is overridden LOCALLY. beamer's default reserves
%  1.8em for the mark; with no mark that is 1.8em of dead indent on every
%  line. Ordinary footnotes elsewhere in the deck are untouched.
%
%  The four \setbeamerfont lines are not optional. \fullcite goes through
%  beamer's bibliography drivers, which apply "bibliography entry ..." fonts --
%  absolute sizes that would override \deckcitefootsize and leave the [n]
%  label tiny next to a much larger reference. Unstarred \setbeamerfont
%  changes only the size and leaves series and shape alone, and all four are
%  local to the group, so the references slide at the end is unaffected.
\newcommand{\deck@citefootline}[1]{%
  \begingroup
    \setbeamerfont{bibliography entry author}  {size=\deck@citefootsize}%
    \setbeamerfont{bibliography entry title}   {size=\deck@citefootsize}%
    \setbeamerfont{bibliography entry location}{size=\deck@citefootsize}%
    \setbeamerfont{bibliography entry note}    {size=\deck@citefootsize}%
    \setbeamertemplate{footnote}{%
      \parindent0pt\noindent\raggedright
      \insertfootnotetext\par
    }%
    \renewcommand{\thefootnote}{}%
    \footnote{{\deck@citefootsize\deckcitenum{#1}\ \fullcite{#1}}}%
    \addtocounter{footnote}{-1}%
  \endgroup
}

% --- Crowding warning ------------------------------------------------------
%  Counts bottom references on the current slide and says so in the log once
%  the slide carries more than \deckcitefootmax of them. A crushed slide
%  should never be a surprise you only discover in the PDF.
\newcount\deck@citefootn
\def\deck@citefootwhere{}
\newcommand{\deck@citefootcount}{%
  \edef\deck@citefoothere{\the\c@page.\the\beamer@slideinframe}%
  \ifx\deck@citefoothere\deck@citefootwhere
    \global\advance\deck@citefootn\@ne
  \else
    \global\deck@citefootn\@ne
    \global\let\deck@citefootwhere\deck@citefoothere
  \fi
  \ifnum\deck@citefootn>\deck@citefootmax\relax
    \PackageWarning{deck-bib}{%
      \the\deck@citefootn\space bottom references on slide
      \the\c@page\space -- more than \the\deck@citefootmax.\MessageBreak
      They will crowd the slide. Use \string\citequiet\space for the
      extras,\MessageBreak
      or \string\deckcitefootoff\space at the top of that frame%
    }%
  \fi
}

% --- Dedup, per slide ------------------------------------------------------
%  Cite the same work twice on one slide and it should appear once at the
%  bottom; cite it again on a later slide and it should appear there too.
%
%  The marker is keyed on the physical page AND on which overlay of the frame
%  is being built. beamer typesets a frame body once per overlay, so a flag
%  reset by a begin-frame hook would fire once and then suppress the block on
%  overlays 2..n -- the reference would vanish the moment you pressed the
%  clicker. Keying on the slide being built needs no hook at all.
\newcommand{\deck@citefootone}[1]{%
  \ifcsname deck@seen@\the\c@page @\the\beamer@slideinframe @#1\endcsname
  \else
    \expandafter\gdef
      \csname deck@seen@\the\c@page @\the\beamer@slideinframe @#1\endcsname{}%
    \deck@citefootcount
    \deck@citefootline{#1}%
  \fi
}

%  \cite{a,b} is legal, so split the list. \forcsvlist trims stray spaces.
\newcommand{\deck@citefoot}[1]{%
  \ifdeck@citefoot
    \forcsvlist{\deck@citefootone}{#1}%
  \fi
}

% --- \cite does both -------------------------------------------------------
%  Wrapped, not reimplemented, so biblatex keeps full control of what [1]
%  actually looks like.
%
%  Deferred to \AtBeginDocument because hyperref -- which beamer loads for us
%  -- patches \cite at the end of the preamble. Grabbing it any earlier would
%  wrap a version that is about to be replaced.
%  biblatex reads ONE optional argument as the postnote and TWO as prenote +
%  postnote, so the wrapper has to hand back exactly as many as it was given.
%  Empty defaults are used rather than -NoValue- markers: an empty argument is
%  something \detokenize can test for certainty, and there is nothing to pass
%  through a second command that might mistake the marker for real text.
\newcommand{\deck@ifempty}[1]{%
  \if\relax\detokenize{#1}\relax
    \expandafter\@firstoftwo
  \else
    \expandafter\@secondoftwo
  \fi
}
\newcommand{\deck@passcite}[3]{%
  \deck@ifempty{#1}%
    {\deck@ifempty{#2}{\deck@origcite{#3}}{\deck@origcite[][#2]{#3}}}%
    {\deck@ifempty{#2}{\deck@origcite[#1]{#3}}{\deck@origcite[#1][#2]{#3}}}%
}

\NewDocumentCommand{\citequiet}{O{} O{} m}{\deck@passcite{#1}{#2}{#3}}

\AtBeginDocument{%
  \let\deck@origcite\cite
  \RenewDocumentCommand{\cite}{s O{} O{} m}{%
    \deck@passcite{#2}{#3}{#4}%
    \IfBooleanF{#1}{\deck@citefoot{#4}}%
  }%
}

\makeatother

%  ---------------------------------------------------------------------------
%      \deckcitefootsize{\scriptsize}  larger bottom text   (default \tiny)
%      \deckcitefootrule{0pt}          drop the hairline above the block
%      \deckcitefootmax{6}             raise the crowding warning threshold
%      \deckcitefootoff                stop adding bottom references
%      \deckcitefooton                 turn them back on
%
%  \deckcitefootoff and \deckcitefooton also work INSIDE a frame, where they
%  apply to that frame only -- useful on a related-work slide with six cites.
%
%  ---------------------------------------------------------------------------
%  WHERE NOT TO PUT \cite
%  ---------------------------------------------------------------------------
%  Not in \frametitle, \section or \caption. Those are moving arguments: the
%  text is written out to the .toc/.aux and typeset again elsewhere, and a
%  footnote cannot survive that. Use \citequiet there.
% ============================================================================

\usepackage[
  backend     = biber,
  style       = numeric,   % [1] in the text, [1] in the reference list
  sorting     = none,      % list ordered by first citation, so [1] really is
                           % the first work you cite -- numbers read in order
  maxbibnames = 6,
  doi         = false,     % slides rarely need these; they eat space
  isbn        = false,
  url         = false,
  eprint      = false,
]{biblatex}

\usepackage{etoolbox}      % \forcsvlist -- biblatex loads it too, this line
                           % is here so the file does not rely on that

% ---------------------------------------------------------------------------
%  THE REFERENCES SLIDE
% ---------------------------------------------------------------------------
%  beamer replaces biblatex's entry label with its own "bibliography item"
%  template. Left empty it drops beamer's generic article icon -- but it also
%  silently deletes the [1] that numeric style depends on. \insertbiblabel
%  puts the number back.
\setbeamertemplate{bibliography item}{\insertbiblabel}
\renewcommand*{\bibfont}{\scriptsize}

% Keep entries from being hyphenated into a mess in a narrow column
\appto{\bibsetup}{\raggedright}

%  A frame that breaks across pages is titled "References I", "References II"
%  and so on. Numbering a list that turned out to fit on one slide just looks
%  like a mistake, so the numeral starts from the second page.
\setbeamertemplate{frametitle continuation}[from second]

% ---------------------------------------------------------------------------
%  A references slide that splits across pages automatically if it is long.
%  Use \bibliographyframe in your deck instead of writing this out by hand.
% ---------------------------------------------------------------------------
\newcommand{\bibliographyframe}[1][References]{%
  \begin{frame}[allowframebreaks]{#1}
    \printbibliography[heading=none]
  \end{frame}%
}

% ===========================================================================
%  THE PER-SLIDE REFERENCE BLOCK
% ===========================================================================
\makeatletter

% --- Knobs -----------------------------------------------------------------
\newcommand*{\deck@citefootsize}{\tiny}
\newcommand*{\deckcitefootsize}[1]{\renewcommand*{\deck@citefootsize}{#1}}

\newlength{\deck@citefootrule}
\setlength{\deck@citefootrule}{0.4pt}
\newcommand*{\deckcitefootrule}[1]{\setlength{\deck@citefootrule}{#1}}

\newcount\deck@citefootmax  \deck@citefootmax=4
\newcommand*{\deckcitefootmax}[1]{\deck@citefootmax=#1\relax}

\newif\ifdeck@citefoot  \deck@citefoottrue
\newcommand*{\deckcitefooton}{\deck@citefoottrue}
\newcommand*{\deckcitefootoff}{\deck@citefootfalse}

% --- The number, on its own ------------------------------------------------
%  Prints just [1] for a key, so the bottom line can be labelled with the
%  same number the reader just saw in the body text.
\DeclareCiteCommand{\deckcitenum}[\mkbibbrackets]
  {}
  {\printfield{labelnumber}}
  {\multicitedelim}
  {}

% --- The hairline above the block ------------------------------------------
%  This is LaTeX's footnote rule, so an ordinary \footnote gets the same line.
%  \deckcitefootrule{0pt} makes it disappear without disturbing the spacing.
\renewcommand{\footnoterule}{%
  \kern-3pt%
  \hbox{\color{BrandRule}%
        \vrule width 0.35\textwidth height \deck@citefootrule depth 0pt}%
  \kern2.6pt%
}

% --- One bottom line -------------------------------------------------------
%  An unmarked footnote: beamer already places footnotes at the bottom of the
%  frame, above the footer, and keeps them working under [allowframebreaks].
%  The mark is emptied and the footnote counter rolled back, because the line
%  carries its own [n] instead.
%
%  The footnote template is overridden LOCALLY. beamer's default reserves
%  1.8em for the mark; with no mark that is 1.8em of dead indent on every
%  line. Ordinary footnotes elsewhere in the deck are untouched.
%
%  The four \setbeamerfont lines are not optional. \fullcite goes through
%  beamer's bibliography drivers, which apply "bibliography entry ..." fonts --
%  absolute sizes that would override \deckcitefootsize and leave the [n]
%  label tiny next to a much larger reference. Unstarred \setbeamerfont
%  changes only the size and leaves series and shape alone, and all four are
%  local to the group, so the references slide at the end is unaffected.
\newcommand{\deck@citefootline}[1]{%
  \begingroup
    \setbeamerfont{bibliography entry author}  {size=\deck@citefootsize}%
    \setbeamerfont{bibliography entry title}   {size=\deck@citefootsize}%
    \setbeamerfont{bibliography entry location}{size=\deck@citefootsize}%
    \setbeamerfont{bibliography entry note}    {size=\deck@citefootsize}%
    \setbeamertemplate{footnote}{%
      \parindent0pt\noindent\raggedright
      \insertfootnotetext\par
    }%
    \renewcommand{\thefootnote}{}%
    \footnote{{\deck@citefootsize\deckcitenum{#1}\ \fullcite{#1}}}%
    \addtocounter{footnote}{-1}%
  \endgroup
}

% --- Crowding warning ------------------------------------------------------
%  Counts bottom references on the current slide and says so in the log once
%  the slide carries more than \deckcitefootmax of them. A crushed slide
%  should never be a surprise you only discover in the PDF.
\newcount\deck@citefootn
\def\deck@citefootwhere{}
\newcommand{\deck@citefootcount}{%
  \edef\deck@citefoothere{\the\c@page.\the\beamer@slideinframe}%
  \ifx\deck@citefoothere\deck@citefootwhere
    \global\advance\deck@citefootn\@ne
  \else
    \global\deck@citefootn\@ne
    \global\let\deck@citefootwhere\deck@citefoothere
  \fi
  \ifnum\deck@citefootn>\deck@citefootmax\relax
    \PackageWarning{deck-bib}{%
      \the\deck@citefootn\space bottom references on slide
      \the\c@page\space -- more than \the\deck@citefootmax.\MessageBreak
      They will crowd the slide. Use \string\citequiet\space for the
      extras,\MessageBreak
      or \string\deckcitefootoff\space at the top of that frame%
    }%
  \fi
}

% --- Dedup, per slide ------------------------------------------------------
%  Cite the same work twice on one slide and it should appear once at the
%  bottom; cite it again on a later slide and it should appear there too.
%
%  The marker is keyed on the physical page AND on which overlay of the frame
%  is being built. beamer typesets a frame body once per overlay, so a flag
%  reset by a begin-frame hook would fire once and then suppress the block on
%  overlays 2..n -- the reference would vanish the moment you pressed the
%  clicker. Keying on the slide being built needs no hook at all.
\newcommand{\deck@citefootone}[1]{%
  \ifcsname deck@seen@\the\c@page @\the\beamer@slideinframe @#1\endcsname
  \else
    \expandafter\gdef
      \csname deck@seen@\the\c@page @\the\beamer@slideinframe @#1\endcsname{}%
    \deck@citefootcount
    \deck@citefootline{#1}%
  \fi
}

%  \cite{a,b} is legal, so split the list. \forcsvlist trims stray spaces.
\newcommand{\deck@citefoot}[1]{%
  \ifdeck@citefoot
    \forcsvlist{\deck@citefootone}{#1}%
  \fi
}

% --- \cite does both -------------------------------------------------------
%  Wrapped, not reimplemented, so biblatex keeps full control of what [1]
%  actually looks like.
%
%  Deferred to \AtBeginDocument because hyperref -- which beamer loads for us
%  -- patches \cite at the end of the preamble. Grabbing it any earlier would
%  wrap a version that is about to be replaced.
%  biblatex reads ONE optional argument as the postnote and TWO as prenote +
%  postnote, so the wrapper has to hand back exactly as many as it was given.
%  Empty defaults are used rather than -NoValue- markers: an empty argument is
%  something \detokenize can test for certainty, and there is nothing to pass
%  through a second command that might mistake the marker for real text.
\newcommand{\deck@ifempty}[1]{%
  \if\relax\detokenize{#1}\relax
    \expandafter\@firstoftwo
  \else
    \expandafter\@secondoftwo
  \fi
}
\newcommand{\deck@passcite}[3]{%
  \deck@ifempty{#1}%
    {\deck@ifempty{#2}{\deck@origcite{#3}}{\deck@origcite[][#2]{#3}}}%
    {\deck@ifempty{#2}{\deck@origcite[#1]{#3}}{\deck@origcite[#1][#2]{#3}}}%
}

\NewDocumentCommand{\citequiet}{O{} O{} m}{\deck@passcite{#1}{#2}{#3}}

\AtBeginDocument{%
  \let\deck@origcite\cite
  \RenewDocumentCommand{\cite}{s O{} O{} m}{%
    \deck@passcite{#2}{#3}{#4}%
    \IfBooleanF{#1}{\deck@citefoot{#4}}%
  }%
}

\makeatother
